% Fragmento de inserción; se incluye desde la raíz de este paquete editorial.
% Las fuentes completas y los cortes materiales constan en ANTECEDENTES_Y_DEPENDENCIAS.md.
% Procedencia: XI REV11 ES, sections/nucleo.tex, tpk_desarrollo_integrado.tex,
% ch_tres_vueltas_ext.tex, ch_elevacion_catalogal.tex, ch_cierre_cardinal.tex,
% ch_descenso_por_fibras.tex; perfil integral public_final/feigenbaum/c37_feigenbaum.tex;
% Grassmannianos REV07 ES, sections/jacobi_spectral.tex;
% TRANSVERSALIDAD_Y_ESCALADO_LOCAL_FEIGENBAUM.md, sección 14.
% Ampliación expositiva de demostraciones existentes; no introduce resultados nuevos.
\providecommand{\ZZ}{\mathbb Z}
\providecommand{\RR}{\mathbb R}
\providecommand{\FF}{\mathbb F}
\providecommand{\APP}{\mathrm{APP}}
\providecommand{\TPK}{\mathrm{TPK}}
\providecommand{\TRIT}{\{-1,0,+1\}}

\section{Antécédents operatoriels de la famille de criticité}
\label{hmtfg:antecedentes-operatorios}

La famille de criticité est construite à partir d’une lecture du calendrier TPK. Les deux feuilles arithmétiques, l’orientation tritique et le transport avec mémoire déterminent d’abord les états et leurs prolongements. La lecture ordonnée du continu permet ensuite d’évaluer des fonctions analytiques, de comparer des paramètres et de sélectionner des racines. Ce développement conserve l’état généalogique et ses deux lectures : la carte archimédienne fournit l’espace d’évaluation, tandis que la lecture d’incidence conserve les relations de frontière. La structure discrète conjointe du continu maintient ses cinq constructions consubstantielles. La présente exposition utilise son lecteur ordonné et conserve l’antécédent commun dont il provient.

\subsection{Cellules arithmétiques, orientation et transport}
\label{hmtfg:app-trit-transporte}

On utilise la cellule relevée \eqref{eq:app-levantada}, les coordonnées tritiques et leur loi arithmétique (\ref{trit-ampl-def-coordenadas}, \ref{trit-ampl-prop-aritmetica}), ainsi que le sélecteur~\ref{tpk-int:seleccion}. Fixons les notations : $D_9=\{1,\ldots,9\}$, $X_9=(\mathbb Z/9\mathbb Z)^2$ et, pour tout $n\in\mathbb Z$, $\rho_9(n)=1+((n-1)\bmod9)$ et $q_9(n)=(n-\rho_9(n))/9$. Le représentant équilibré est $r_3(n)\in\{-1,0,+1\}$ ; sa retenue est $\chi(a,b)=(a+b-r_3(a+b))/3$. Ces formules conservent l’extension aux entiers des évaluations positives du noyau.

Le transport relevé et sa réversibilité sont démontrés en \ref{tpk-int:prop-desplazamiento} ; la configuration observable et son retour de $108$ pas, en \ref{tpk-int:fobs} et \ref{tpk-int:periodo}. L’actualisation \eqref{eq:actualizacion} conserve l’état complet. L’holonomie de \ref{tpk-int:gamma} et les retours de \ref{hist:tres-retornos} distinguent restauration de phase et progression de la mémoire. La réalisation concrète des trois branches et son registre intégral sont construits ci-dessous à partir de ces opérations.

% Procedencia literal: output/PREPARACION_ENTREGA_USB_20260928/ENTREGA_HMT_USB/03_FUENTES_Y_REPRODUCCION/ACTUALIZACIONES_20260928/FINAL01/PAQUETES/ARTICULO_XI_REV11_ES_EN_FINAL_BN_20260928/ARTICULO_XI_REV11_ES_EN_FINAL_BN/ES/source/sections/ch_tres_vueltas_ext.tex
% Recuperación íntegra; única transformación mecánica: prefijo hmtfg: en label/ref/eqref.
% No se modifican operadores, pruebas, tablas ni afirmaciones del propietario.
\subsubsection{Réalisation arête par arête des trois prolongements nonadiques}
\label{hmtfg:subsubsec:tres-vueltas-ext-ch}

Cette sous-section spécifie le modèle autonome engendré par la relation
d’extension TPK qui intervient dans la preuve cardinale. On n’infère pas l’existence
d’une arête par coïncidence de types : on construit ses états initial et final,
sa relation APP, son transport de fibre et la mise à jour unaire du registre.
La construction utilise exclusivement les deux feuillets d’APP, le sélecteur TRIT
et la composition \(\mathsf{Sel}\)--\(\mathsf{Tra}\)--\(\operatorname{Upd}\) du TPK.

\paragraph{Cellules complètes et trois réductions équilibrées.}
Dans la carte \(D_9=\{1,\ldots,9\}\), nous maintiendrons toujours la cellule APP complète
\[
 a_{p,q}:=(p,q;R^\Sigma_{p,q},q^\Sigma_{p,q},
                 R^\Pi_{p,q},q^\Pi_{p,q}),
\]
avec
\[
 p+q=R^\Sigma_{p,q}+9q^\Sigma_{p,q},\qquad
 pq=R^\Pi_{p,q}+9q^\Pi_{p,q}.
\]
Le feuillet actif est enregistré dans une coordonnée séparée ; \(a_{p,q}\) n'est pas scindé en une prétendue cellule additive et une autre multiplicative. Pour \(r\in D_9\), soient
\begin{equation}
 a(r):=
 \begin{cases}
  a_{1,9},&r=1,\\
  a_{1,r-1},&2\le r\le9,
 \end{cases}
 \qquad
 \tau(r):=M_{\rm ph}(r),
 \qquad
 m(r):=\tau(r)\in\{-1,0,+1\}.
 \label{hmtfg:eq:celula-completa-fase-ch}
\end{equation}
Le reste additif de \(a(r)\) fournit ici un représentant APP de la
phase ; cela ne transforme pas \(\Sigma\) en feuillet actif universel. Le sélecteur
dynamique active \(\Sigma\) dans \(r\in\{1,4,7\}\), active \(\Pi\) dans
\(r\in\{2,5,8\}\) et conserve les deux feuillets, sans déplacer le curseur, dans
\(r\in\{3,6,9\}\). Dans les trois cas, on transporte conjointement
\(R^\Sigma,q^\Sigma,R^\Pi,q^\Pi\). On a
\[
 (m(1),\ldots,m(9))=(1,-1,0,1,-1,0,1,-1,0).
\]
Les trois paires ordonnées par la base nonadique sont
\begin{equation}
 \mathfrak p_{+1}:=(1,4),\qquad
 \mathfrak p_{-1}:=(2,5),\qquad
 \mathfrak p_{0}:=(3,6).
 \label{hmtfg:eq:tres-pares-internos-ch}
\end{equation}
Ces couples constituent la section normalisée qui prend les deux premières occurrences de chaque type TRIT sous le décalage \(r\mapsto r+3\) ; le troisième bloc de phases \(7,8,9\) conserve le contrôle de fermeture du tour. Leurs deux membres appartiennent à une même fibre de \(M_{\rm ph}\). La réduction équilibrée définie dans ces antécédents donne, sans introduire de coefficient cible,
\begin{equation}
 \begin{array}{c|c|c|c}
 \varepsilon&m(\mathfrak p_\varepsilon^{(1)})
                  +m(\mathfrak p_\varepsilon^{(2)})
   &r_3\bigl(2\varepsilon\bigr)&
   \chi(\varepsilon,\varepsilon)\\ \hline
 +1&1+1=(-1)+3\cdot1&-1&+1\\
 0&0+0=0+3\cdot0&0&0\\
 -1&(-1)+(-1)=1+3\cdot(-1)&+1&-1
 \end{array}
 \label{hmtfg:eq:tres-carry-derivados-ch}
\end{equation}
Par conséquent,
\begin{equation}
 \chi(\varepsilon,\varepsilon)
 =\frac{2\varepsilon-r_3(2\varepsilon)}3=\varepsilon.
 \label{hmtfg:eq:carry-no-parametrico-ch}
\end{equation}
Le signe de la branche est ainsi la sortie du cocycle équilibré appliqué à deux
émissions TRIT produites par APP. Chaque paire fournit l’une des trois
sorties de retenue. Par souci de clarté, nous indexerons ensuite la paire par
\(\varepsilon(\mathfrak p):=
\chi(m(\mathfrak p^{(1)}),m(\mathfrak p^{(2)}))\) ; le symbole
\(\varepsilon\) ne fournit pas à \(\operatorname{Upd}^{\rm cal}\) une retenue libre.

L’architecture de phase, le relèvement équilibré et le transport de
mémoire sont réalisés ici au moyen de vingt-sept incidences locales et de leurs
graphes d’extension. La section normalisée fixe des représentants des
fibres de \(M_{\rm ph}\) ; la couverture qui sera démontrée utilise l’existence
de ces représentants et demeure invariante lorsqu’on les remplace par une autre
section compatible.

\paragraph{Données typées des neuf arêtes.}
Soit
\(K_{\rm Carr}:=\langle\kappa_{\rm Carr}\rangle_{\rm Ab}\) le module de
retenue. Les actions sont typées par
\[
 H_{\rm Carr}:\mathscr P_{\rm TPK}^{\rm op}
   \longrightarrow\operatorname{End}(K_{\rm Carr}),\qquad
 Y_{\rm Carr}:\mathscr P_{\rm TPK}
   \longrightarrow\operatorname{End}(K_{\rm Carr}).
\]
Dans ce modèle engendré, on prend les actions triviales \(H_{\rm Carr}(e)=Y_{\rm Carr}(e)=\operatorname{id}_{K_{\rm Carr}}\) : le feuillet actif et le chemin sont conservés comme coordonnées propres du registre et n'agissent pas sur le terme libre de retenue. Soit en outre \((\mathsf M_d,\jmath_{d',d})\) le système inductif gradué de mémoire. Pour ne pas confondre deux représentants identifiés par la colimite, l'état conserve en outre le degré comme coordonnée et l'on utilise l'espace porteur étiqueté
\[
 \widetilde{\mathsf M}:=
 \bigsqcup_{d\ge1}\bigl(\{d\}\times\mathsf M_d\bigr),
 \qquad
 \operatorname{gr}_{\mathsf M}(d,\mu)=d.
\]
La publication dans la mémoire limite est donnée par
\[
\begin{aligned}
 \mathsf M&:=\varinjlim_d\mathsf M_d,&
 \jmath_d&:\mathsf M_d\longrightarrow\mathsf M,\\
 u&:=\jmath_d(u_d),&
 \iota_{\mathsf M}:\mathbb Z&\longrightarrow\mathsf M,
 \qquad n\longmapsto nu .
\end{aligned}
\]
Ici, \(\jmath_d\) est l’application canonique de la colimite et les \(u_d\) forment une
famille compatible. Sa composante libre permet
le lecteur \(\deg_u(nu)=n\).
Les projections de configuration, de phase, de mémoire, de retenue, de route et de
survie seront typées sur le porteur calibré construit
récursivement ci-dessous ; aucun espace ambiant ne sera importé pour elles.
Pour chaque arête construite \(e\), la
composante calibrée
\(\operatorname{Upd}^{\rm cal}_e:=\mathsf{Mem}^{\rm cal}_e\) désignera la
troisième application de la factorisation ; sa loi est fixée champ par champ plus
bas et n’est pas présupposée comme restriction d’une relation ambiante.

Fixons la profondeur initiale \(k_0=1\), une configuration complète
\(q_\star\in\mathcal Q_{\rm obs}\) de phase un et
\[
 q_{n,r}:=F_{\rm obs}^{9n+r-1}(q_\star),\qquad
 q_{n,10}:=q_{n+1,1}.
\]
Comme le calendrier observable est périodique, sa valeur \(q_{n,r}\) ne
distingue pas à elle seule des apparitions séparées par douze tours. Le modèle
calibré emploie donc le relèvement étiqueté
\[
 \widetilde q_{n,r}:=(q_{n,r},n,r),\qquad
 \widetilde q_{n,10}:=\widetilde q_{n+1,1},\qquad
 \pi_{\rm obs}(\widetilde q_{n,r})=q_{n,r}.
\]
La dynamique élevée \(\widetilde F_{\rm obs}(\widetilde q_{n,r})
=\widetilde q_{n,r+1}\) se projette sur \(F_{\rm obs}\) ; l'étiquette ne modifie pas la phase physique, mais distingue ses niveaux d'apparition. On étiquette la copie de la cellule à chaque niveau au moyen de
\[
 a_{n,r}:=(n,r;a(r)),\qquad
 (n,r)^+:=
 \begin{cases}
  (n,r+1),&1\le r<9,\\
  (n+1,1),&r=9.
 \end{cases}
\]
L’étiquette évite seulement d’identifier deux apparitions d’une même cellule ; les
deux divisions euclidiennes et leurs quatre coordonnées sont celles de \(a(r)\).
Pour le tour \(n\ge0\) et la phase \(r\in D_9\), le symbole
\[
 e_{\varepsilon,n,r}:\widetilde q_{n,r}
 \longrightarrow\widetilde q_{n,r+1},
 \qquad r^+:=1+(r\bmod9),
\]
occupe le niveau \(d_{n,r}:=k_n+r-1\), où \(k_n:=k_0+9n\). Son registre local contient la cellule complète \(a(r)\), le feuillet actif déterminé par \(M_{\rm ph}(r)\), le type \(\tau(r)\), la marque d'appartenance à \(\mathfrak p_\varepsilon\) et, lorsqu'il atteint le second membre de cette paire, l'identité certifiée \eqref{hmtfg:eq:tres-carry-derivados-ch}. Nous définissons ses opérateurs uniquement à partir de ce registre :
\begin{align}
 \mathsf C_{\mathsf M}(e_{\varepsilon,n,r})
   &=\jmath_{d_{n,r}+1,d_{n,r}},&
 \kappa_{\mathsf M}(e_{\varepsilon,n,r})
   &=\mathbf1_{\{r=9\}}u_{d_{n,r}+1},
 \label{hmtfg:eq:memoria-arista-ch}\\
       \operatorname{Carr}(e_{\varepsilon,n,r})
   &=d_{\varepsilon,r}\kappa_{\rm Carr},&
 d_{\varepsilon,r}
   &:=\mathbf1_{\{r=\mathfrak p_\varepsilon^{(2)}\}}
       \chi\!\left(m(\mathfrak p_\varepsilon^{(1)}),
                    m(\mathfrak p_\varepsilon^{(2)})\right).
 \label{hmtfg:eq:carry-arista-derivado-ch}
\end{align}
Pour les identités et pour toute composition admissible, on définit en outre
\begin{align}
 \mathsf C_{\mathsf M}(\operatorname{id}_{\widetilde q};d)
   &:=\operatorname{id}_{\mathsf M_d},&
 \kappa_{\mathsf M}(\operatorname{id}_{\widetilde q};d)&:=0,
 \label{hmtfg:eq:memoria-identidad-ch}\\
 \mathsf C_{\mathsf M}(e_t\cdots e_1)
   &:=\mathsf C_{\mathsf M}(e_t)\circ\cdots\circ
      \mathsf C_{\mathsf M}(e_1),&
 \kappa_{\mathsf M}(e_2e_1)
   &:=\mathsf C_{\mathsf M}(e_2)
        \bigl(\kappa_{\mathsf M}(e_1)\bigr)
      +\kappa_{\mathsf M}(e_2).
 \label{hmtfg:eq:cociclo-memoria-calibrado-ch}
\end{align}
L’identité porte le degré \(d\) car une même configuration observable
peut réapparaître à des profondeurs distinctes. La dernière égalité, itérée, définit
\(\kappa_{\mathsf M}(e_t\cdots e_1)\) pour toute longueur et est la loi de
cocycle du système direct ; les inclusions satisfont
\(\jmath_{d+2,d+1}\jmath_{d+1,d}=\jmath_{d+2,d}\).
On fixe en outre
\(s(e_{\varepsilon,n,r})=\top\). Les vingt-sept entrées suivantes
certifient les coordonnées qui distinguent les trois prolongements ; les champs
universels du transport sont spécifiés immédiatement après. Les symboles \(I\) et \(II\)
désignent respectivement la première émission conservée et la seconde émission
qui effectue la réduction ; un tiret indique un accroissement de retenue nul, mais non
une absence d’arête.
\begin{center}
\small
\setlength{\tabcolsep}{3.2pt}
\begin{tabular}{c c cc cc cc c}
\toprule
\(r\)&\(M_{\rm ph}(r)\)&\multicolumn{2}{c}{\(\mathfrak p_{+1}\)}&
\multicolumn{2}{c}{\(\mathfrak p_0\)}&
\multicolumn{2}{c}{\(\mathfrak p_{-1}\)}&
\(\Delta\operatorname{mem}\)\\
& &marque&\(d_{+1,r}\)&marque&\(d_{0,r}\)&marque&\(d_{-1,r}\)&\\
\midrule
1&\(+1\)&I&0&--&0&--&0&0\\
2&\(-1\)&--&0&--&0&I&0&0\\
3&0&--&0&I&0&--&0&0\\
4&\(+1\)&II&\(+1\)&--&0&--&0&0\\
5&\(-1\)&--&0&--&0&II&\(-1\)&0\\
6&0&--&0&II&0&--&0&0\\
7&\(+1\)&--&0&--&0&--&0&0\\
8&\(-1\)&--&0&--&0&--&0&0\\
9&0&--&0&--&0&--&0&\(u\)\\
\bottomrule
\end{tabular}
\end{center}
Ainsi, \(d_{\varepsilon,r}\) est calculé à partir de deux sorties du sélecteur et
est distinct des quotients \(q^\Sigma,q^\Pi\) de la cellule. Le registre
conserve également la paire marquée et le résidu équilibré ; en particulier, la
branche nulle n’est pas confondue avec une absence d’événement.

La loi de composition de la retenue employée par la catégorie TPK est
\begin{equation}
 \operatorname{Carr}(\operatorname{id}_{\widetilde q})=0,\qquad
 \operatorname{Carr}(e_2e_1)
 =H_{\rm Carr}(e_1)\bigl(\operatorname{Carr}(e_2)\bigr)
  +Y_{\rm Carr}(e_2)\bigl(\operatorname{Carr}(e_1)\bigr).
 \label{hmtfg:eq:ley-carry-correcta-ch}
\end{equation}
Sur les identités et les chemins, ils se prolongent fonctoriellement :
\[
\begin{aligned}
 H_{\rm Carr}(\operatorname{id})
   &=Y_{\rm Carr}(\operatorname{id})=\operatorname{id},\\
 H_{\rm Carr}(e_t\cdots e_1)
   &=H_{\rm Carr}(e_1)\circ\cdots\circ H_{\rm Carr}(e_t),\\
 Y_{\rm Carr}(e_t\cdots e_1)
   &=Y_{\rm Carr}(e_t)\circ\cdots\circ Y_{\rm Carr}(e_1).
\end{aligned}
\]
Dans ce modèle, toutes ces actions restent l'identité, mais leur typage explicite quel endomorphisme agit sur chaque terme de \eqref{hmtfg:eq:ley-carry-correcta-ch}.
Sur les arêtes calibrées, \(H_{\rm Carr}=Y_{\rm Carr}=\operatorname{id}\), si bien que la composition additionne les accroissements déjà produits. L'équation \eqref{hmtfg:eq:ley-carry-correcta-ch}, et non une identification entre orientation et quotient, gouverne le registre accumulé.

\paragraph{Construction indépendante des tuples de transition.}
Un état de la fibre sera écrit, dans l’ordre de ses champs pertinents,
\[
 x=(a;q;\tau,o,h;b_{\rm ref};
       R^\Sigma,q^\Sigma,R^\Pi,q^\Pi;c;
       \operatorname{Sig};C,\partial C;p_{\APP},p_{\TRIT};
       \lambda;\operatorname{gr}_{\mathsf M};\operatorname{mem};
       s;\operatorname{Led}).
\]
La notation de mise à jour \(x[f\leftarrow v]\) ne change que le champ
indiqué. Pour \(h\in\{(\Sigma,\Pi),(\Pi,\Sigma)\}\), définissons
\[
 h_r(h)=
 \begin{cases}
  (\Sigma,\Pi),&r\in\{1,4,7\},\\
  (\Pi,\Sigma),&r\in\{2,5,8\},\\
 h,&r\in\{3,6,9\}.
 \end{cases}
\]
La coordonnée complète de feuillet et d’arithmétique est abrégée par
\[
 \widehat h(x):=
 \bigl(h(x),R^\Sigma(x),q^\Sigma(x),R^\Pi(x),q^\Pi(x)\bigr).
\]
Soit \(\widehat{\mathcal H}^{\rm cal}_{\widetilde q_{n,r}}\) l’ensemble des
quintuplets
\[
 \bigl(h,R^\Sigma_{a_{n,r}},q^\Sigma_{a_{n,r}},
          R^\Pi_{a_{n,r}},q^\Pi_{a_{n,r}}\bigr),
 \qquad
 h\in\{(\Sigma,\Pi),(\Pi,\Sigma)\},
\]
dont les quatre données arithmétiques sont les deux divisions euclidiennes de la cellule
APP étiquetée \(a_{n,r}\). Cette définition directe ne présuppose pas d’état
sur un porteur qui n’est pas encore construit. Soient
\(\mathsf{Hist}^{\rm cal}_{\widetilde q}\) l’ensemble des chemins finis
admissibles du graphe d’arêtes étiquetées qui aboutissent à
\(\widetilde q\), y compris l’identité, et
\(\mathsf B^{\rm cal}_{\widetilde q}\) l’ensemble de leurs suites
ordonnées d’extrémités. Nous définissons
\[
 \mathsf S^{\rm cal}_{\widetilde q}:=
 \widehat{\mathcal H}^{\rm cal}_{\widetilde q}
 \times\{o_{\rightarrow},o_{\circ},o_{\leftarrow}\}
 \times K_{\rm Carr}\times
 \mathsf{Hist}^{\rm cal}_{\widetilde q}\times
 \mathsf B^{\rm cal}_{\widetilde q}\times\mathsf M .
\]
La signature calibrée n’est pas postulée : c’est l’application de réunion typée
\[\begin{gathered}
 \operatorname{Sig}^{\rm cal}_{\widetilde q}:
 \widehat{\mathcal H}^{\rm cal}_{\widetilde q}
 \times\{o_{\rightarrow},o_{\circ},o_{\leftarrow}\}
 \times K_{\rm Carr}\times
 \mathsf{Hist}^{\rm cal}_{\widetilde q}\times
 \mathsf B^{\rm cal}_{\widetilde q}\times\mathsf M
 \longrightarrow\mathsf S^{\rm cal}_{\widetilde q},
\\
 (\widehat h,o,c,\lambda,\partial C,\mu)
 \longmapsto(\widehat h,o,c,\lambda,\partial C,\mu).
\end{gathered}\]
De même, \(\pi_{\mathcal F}\) désigne la projection qui élimine uniquement les deux premières coordonnées \((a;q)\) du tuple d'état ; elle conserve \(\tau,o,h,b_{\rm ref}\), les quatre données euclidiennes, la retenue, la signature, le cylindre, la frontière, les préfixes, le chemin, le degré de mémoire, la mémoire, la survie et le registre.
L'orientation correspondante est fixée par \(o(+1)=o_{\rightarrow}\), \(o(0)=o_{\circ}\) et \(o(-1)=o_{\leftarrow}\). La fibre APP calibrée sur \(\widetilde q_{n,r}\) contiendra la copie étiquetée \(a_{n,r}\) déjà définie. L'état initial est fixé sans coordonnées libres :
\[
 \widehat h_\star:=
 \bigl((\Sigma,\Pi),R^\Sigma_{1,9},q^\Sigma_{1,9},
       R^\Pi_{1,9},q^\Pi_{1,9}\bigr).
\]
\begin{equation}
\begin{aligned}
 x_\star:=\bigl(&a_{0,1};\widetilde q_{0,1};
 +1,o_{\rightarrow},(\Sigma,\Pi);\bot;
 R^\Sigma_{1,9},q^\Sigma_{1,9},R^\Pi_{1,9},q^\Pi_{1,9};0;\\
 &\operatorname{Sig}^{\rm cal}_{\widetilde q_{0,1}}
       (\widehat h_\star,o_{\rightarrow},0,
        \operatorname{id}_{\widetilde q_{0,1}},\varnothing,
        \jmath_{k_0}(0_{\mathsf M_{k_0}}));
 \varnothing,\varnothing;
 (a_{0,1}),(+1);\operatorname{id}_{\widetilde q_{0,1}};
 k_0;0_{\mathsf M_{k_0}};\top;\varnothing\bigr).
\end{aligned}
\label{hmtfg:eq:estado-inicial-calibrado-ch}
\end{equation}
Le tuple ne contient pas de coordonnées indéterminées et satisfait
\(s(x_\star)=\top\). Soit \(G_0:=\{x_\star\}\).
Supposant \(G_n\) construit, fixons \(x\in G_n\),
\(\varepsilon\in\{-1,0,+1\}\) et
\(x_{\varepsilon,0}(x):=x\). Pour \(r=1,\ldots,9\), on construit d’abord
le registre élémentaire
\begin{equation}
\begin{aligned}
 \ell_{\varepsilon,n,r}:=\bigl(&e_{\varepsilon,n,r},a_{n,r},a_{(n,r)^+},
 \widetilde q_{n,r},\widetilde q_{n,r+1},M_{\rm ph}(r),\\
 &h_r(h(x_{\varepsilon,r-1}(x))),o(M_{\rm ph}(r)),\mathfrak p_\varepsilon,
 d_{\varepsilon,r},\\
 &\mathsf C_{\mathsf M}(e_{\varepsilon,n,r}),
 \kappa_{\mathsf M}(e_{\varepsilon,n,r}),
 \operatorname{Carr}(e_{\varepsilon,n,r})\bigr),
\end{aligned}
 \label{hmtfg:eq:registro-elemental-ext-ch}
\end{equation}
et on pose
\(\operatorname{Led}^{\rm new}_{\varepsilon,n,r}:=
\operatorname{Led}(x_{\varepsilon,r-1}(x))^\frown
\ell_{\varepsilon,n,r}\). On construit alors le tuple
\begin{equation}
 y_{\varepsilon,r}(x):=
 x_{\varepsilon,r-1}(x)
 [\tau\leftarrow M_{\rm ph}(r),\
  h\leftarrow h_r(h(x_{\varepsilon,r-1}(x))),\
  o\leftarrow o(M_{\rm ph}(r)),\
  b_{\rm ref}\leftarrow\mathfrak p_\varepsilon],
 \label{hmtfg:eq:tupla-sel-cal-ch}
\end{equation}
où la sélection laisse intacts la cellule, les quotients, la retenue, la route, la mémoire,
le cylindre, la frontière, la survie et le registre ; la coordonnée
\(b_{\rm ref}\) explicite la réduction équilibrée choisie. On ne présuppose pas ici
l’appartenance à une relation ambiante.

On définit ensuite \(z_{\varepsilon,r}(x)\) champ par champ : sa
configuration est \(\widetilde q_{n,r+1}\) ; sa cellule complète est
\(a_{(n,r)^+}\), avec les
quatre résidus et quotients recalculés par les deux divisions euclidiennes ;
il conserve le feuillet sélectionné et l’orientation ; et il effectue les six
mises à jour suivantes, dans lesquelles seul l’argument commun \(x\) est omis :
\begin{equation}
\begin{aligned}
 C(z_{\varepsilon,r})&:=C(y_{\varepsilon,r})^\frown
       e_{\varepsilon,n,r},\\
 \partial C(z_{\varepsilon,r})&:=\partial C(y_{\varepsilon,r})^\frown
       (\widetilde q_{n,r},\widetilde q_{n,r+1}),\\
 p_{\APP}(z_{\varepsilon,r})&:=p_{\APP}(y_{\varepsilon,r})^\frown
       a_{(n,r)^+},\\
 p_{\TRIT}(z_{\varepsilon,r})&:=p_{\TRIT}(y_{\varepsilon,r})^\frown
       \tau(r),\\
 \operatorname{Led}(z_{\varepsilon,r})&:=
       \operatorname{Led}^{\rm new}_{\varepsilon,n,r},\\
 \lambda(z_{\varepsilon,r})&:=e_{\varepsilon,n,r}
       \lambda(y_{\varepsilon,r}),
\end{aligned}
 \label{hmtfg:eq:tupla-tra-cal-ch}
\end{equation}
et maintient encore les champs \(c,\operatorname{gr}_{\mathsf M},\operatorname{mem},s\) de \(y_{\varepsilon,r}(x)\). Sa signature, en revanche, est à nouveau typée dans la configuration de destination au moyen de la présignature transportée
\begin{equation}
\begin{aligned}
 \operatorname{Sig}(z_{\varepsilon,r}(x))
  :=\operatorname{Sig}^{\rm cal}_{\widetilde q_{n,r+1}}
       \Bigl(&\widehat h(z_{\varepsilon,r}(x)),
             o(z_{\varepsilon,r}(x)),c(y_{\varepsilon,r}(x)),
             \lambda(z_{\varepsilon,r}(x)),\\[-2pt]
             &\partial C(z_{\varepsilon,r}(x)),
             \jmath_{d_{n,r}}
                (\operatorname{mem}_{d_{n,r}}
                   (y_{\varepsilon,r}(x)))\Bigr).
\end{aligned}
\label{hmtfg:eq:prefirma-transporte-cal-ch}
\end{equation}
Ainsi, et seulement dans
\(\mathsf{Tra}^{\rm cal}\), on ajoute l’arête, ses extrémités, la marque
\((\mathfrak p_\varepsilon,r,d_{\varepsilon,r})\) et ses cocycles au registre des
incidences.

L'état initial satisfait
\[
 c(x_\star)=0=\operatorname{Carr}
   (\operatorname{id}_{\widetilde q_{0,1}})
   =\operatorname{Carr}(\lambda(x_\star)).
\]
Dans les états complets \(x_{\varepsilon,r-1}(x)\), et aussi après la
sélection \(y_{\varepsilon,r}(x)\), on conserve l’invariant typé
\(c(v)=\operatorname{Carr}(\lambda(v))\), puisque
\(\mathsf{Sel}^{\rm cal}\) ne modifie aucun de ces deux champs. Le
transport fait avancer la route mais laisse la retenue en attente ; le tuple de
préactualisation satisfait donc exactement
\[
 c(z_{\varepsilon,r}(x))=c(y_{\varepsilon,r}(x))
 =\operatorname{Carr}(\lambda(y_{\varepsilon,r}(x))).
\]
L'application \(\operatorname{Upd}^{\rm cal}\) restaure alors \(c(x_{\varepsilon,r}(x))=
\operatorname{Carr}(\lambda(x_{\varepsilon,r}(x)))\) au moyen de la loi de composition, sans attribuer cet invariant à l'état intermédiaire \(z\).

Enfin, on applique la fonction unaire
\(\operatorname{Upd}^{\rm cal}_{e_{\varepsilon,n,r}}
=\mathsf{Mem}^{\rm cal}_{e_{\varepsilon,n,r}}\) à ce tuple. Son
résultat \(x_{\varepsilon,r}(x)\) est donné par
\begin{equation}
\begin{aligned}
 c\bigl(x_{\varepsilon,r}(x)\bigr)
   &=\operatorname{Carr}\bigl(\lambda(z_{\varepsilon,r}(x))\bigr)\\
   &=H_{\rm Carr}(\lambda(y_{\varepsilon,r}(x)))
       \bigl(\operatorname{Carr}(e_{\varepsilon,n,r})\bigr)
     +Y_{\rm Carr}(e_{\varepsilon,n,r})
       \bigl(c(y_{\varepsilon,r}(x))\bigr),\\
 \lambda\bigl(x_{\varepsilon,r}(x)\bigr)
   &=\lambda(z_{\varepsilon,r}(x)),\\
 \operatorname{gr}_{\mathsf M}
      \bigl(x_{\varepsilon,r}(x)\bigr)
   &=d_{n,r}+1,\\
 \operatorname{mem}_{d_{n,r}+1}
       \bigl(x_{\varepsilon,r}(x)\bigr)
   &=\mathsf C_{\mathsf M}(e_{\varepsilon,n,r})
       \operatorname{mem}_{d_{n,r}}(z_{\varepsilon,r}(x))
     +\kappa_{\mathsf M}(e_{\varepsilon,n,r}),\\
 s\bigl(x_{\varepsilon,r}(x)\bigr)
   &=s(e_{\varepsilon,n,r})\wedge s(z_{\varepsilon,r}(x)),\\
 \operatorname{Sig}\bigl(x_{\varepsilon,r}(x)\bigr)
   &=\operatorname{Sig}^{\rm cal}_{\widetilde q_{n,r+1}}
       \Bigl(\widehat h(z_{\varepsilon,r}(x)),
             o(z_{\varepsilon,r}(x)),
             c(x_{\varepsilon,r}(x)),
             \lambda(z_{\varepsilon,r}(x)),
             \partial C(z_{\varepsilon,r}(x)),\\[-2pt]
   &\hspace{112pt}
             \jmath_{d_{n,r}+1}
              \bigl(\operatorname{mem}_{d_{n,r}+1}
                 (x_{\varepsilon,r}(x))\bigr)\Bigr).
\end{aligned}
\label{hmtfg:eq:upd-unaria-tres-vueltas-ch}
\end{equation}
Tous les champs non énumérés dans \eqref{hmtfg:eq:upd-unaria-tres-vueltas-ch} coïncident avec ceux de \(z_{\varepsilon,r}(x)\). En particulier, la signature est calculée avec le feuillet, l'orientation, la frontière et le chemin déjà fixés par \(\mathsf{Sel}^{\rm cal}\) et \(\mathsf{Tra}^{\rm cal}\), et avec la retenue et la mémoire calculées dans les lignes précédentes ; elle n'est pas définie circulairement à partir de l'état final. Ni \(\varepsilon\), ni \(d_{\varepsilon,r}\), ni \(u\) ne sont des arguments supplémentaires de \(\operatorname{Upd}^{\rm cal}\) : la fonction lit les trois données de l'arête déjà construite.

\paragraph{Récurrence simultanée des niveaux.}
Une fois les neuf étapes précédentes achevées pour un \(G_n\) déjà construit,
on définit immédiatement
\begin{equation}
 \gamma_{\varepsilon,k_n}(x):=x_{\varepsilon,9}(x),\qquad
 G_{n+1}:=\{\gamma_{\varepsilon,k_n}(x):
       x\in G_n,\ \varepsilon\in\TRIT\}.
 \label{hmtfg:eq:injerto-total-tres-vueltas-ch}
\end{equation}
\label{hmtfg:eq:tres-vueltas-tpk-ch}
La base \(G_0=\{x_\star\}\), les tuples intermédiaires et
\eqref{hmtfg:eq:injerto-total-tres-vueltas-ch} constituent une unique définition par
récurrence sur \(n\). Par conséquent, aucun porteur ultérieur n’est utilisé pour
supposer l’existence de \(G_n\) : le niveau suivant n’est introduit
qu’après avoir construit toutes ses arêtes à partir du niveau précédent.

\paragraph{Relation d’extension TPK et appartenance effective.}
Soit \(\mathcal E^{\rm cal}_{\widetilde q}\) la plus petite famille étiquetée qui contient
\(x_\star\) et, à chaque étape de la récurrence précédente, les quatre tuples
\(x_{\varepsilon,r-1}(x),y_{\varepsilon,r}(x),
z_{\varepsilon,r}(x),x_{\varepsilon,r}(x)\) dans leurs configurations d’origine
et de destination. Soit également
\(\mathcal A^{\rm cal}_{\widetilde q_{n,r}}:=\{a_{n,r}\}\), et posons
\[
 \widetilde{\mathcal Q}_{\rm obs}:=
   \{\widetilde q_{n,r}:n\ge0,\ 1\le r\le9\},
 \qquad
 \mathcal E^{\rm cal}:=\bigsqcup_{\widetilde q}
     \mathcal E^{\rm cal}_{\widetilde q},
 \qquad
 \mathcal F^{\rm cal}_{\widetilde q}:=
     \pi_{\mathcal F}(\mathcal E^{\rm cal}_{\widetilde q}).
\]
Pour chaque profondeur \(d\), soit
\(\mathcal E^{\rm cal}_d:=
\operatorname{gr}_{\mathsf M}^{-1}(d)\subset\mathcal E^{\rm cal}\). La
coordonnée explicite de degré, et non la seule appartenance d’un représentant à
une image du système direct, rend ces porteurs disjoints. Sur eux,
les projections sont typées
\[
\begin{aligned}
 \operatorname{cfg}&:\mathcal E^{\rm cal}\longrightarrow
     \widetilde{\mathcal Q}_{\rm obs},\\
 g(x)&:=\operatorname{fase}\!
       \bigl(\pi_{\rm obs}(\operatorname{cfg}(x))\bigr)\in C_9,\\
 c&:\mathcal E^{\rm cal}\longrightarrow K_{\rm Carr},\\
 s&:\mathcal E^{\rm cal}\longrightarrow\{\bot,\top\},\\
 \lambda&:\mathcal E^{\rm cal}_{\widetilde q}
       \longrightarrow\mathsf{Hist}^{\rm cal}_{\widetilde q},\\
 \operatorname{gr}_{\mathsf M}&:\mathcal E^{\rm cal}
       \longrightarrow\mathbb N_{\ge1},\\
 \operatorname{mem}_d&:\mathcal E^{\rm cal}_d
       \longrightarrow\mathsf M_d,\\
 \operatorname{mem}(x)&:=
       \jmath_d\bigl(\operatorname{mem}_d(x)\bigr)\in\mathsf M
       \quad(x\text{ de profondeur }d).
\end{aligned}
\]
Ici, \(\operatorname{cfg}(x)\) est la seconde coordonnée du tuple et
\(\operatorname{fase}\) lit la phase du calendrier observable. Pour les trois
classes de porteur, on fixe les diagonales
\[
\begin{aligned}
 \Delta_{\APP,\widetilde q}
   &:=\{(a,a):a\in\mathcal A^{\rm cal}_{\widetilde q}\},\\
 \Delta_{{\rm fib},\widetilde q}
   &:=\{(f,f):f\in\mathcal F^{\rm cal}_{\widetilde q}\},\\
 \Delta_{{\rm enr},\widetilde q}
   &:=\{(x,x):x\in\mathcal E^{\rm cal}_{\widetilde q}\}.
\end{aligned}
\]
En particulier,
\[
 a_{0,1}\in\mathcal A^{\rm cal}_{\widetilde q_{0,1}},
 \qquad
 x_\star\in\mathcal E^{\rm cal}_{\widetilde q_{0,1}}.
\]
Sur ces espaces porteurs, on définit par leurs graphes complets, et non au moyen d'implications nécessaires, les relations
\begin{align}
 \mathsf{Sel}^{\rm cal}_{e_{\varepsilon,n,r}}
   &:=\{(x_{\varepsilon,r-1}(x),y_{\varepsilon,r}(x)):
          x\in G_n\},\\
 \mathsf{Tra}^{\rm cal}_{e_{\varepsilon,n,r}}
   &:=\{(y_{\varepsilon,r}(x),z_{\varepsilon,r}(x)):
          x\in G_n\},\\
 \operatorname{Ext}^{\APP,\rm cal}_{e_{\varepsilon,n,r}}
   &:=\{(a_{n,r},a_{(n,r)^+})\},
 \label{hmtfg:eq:ext-app-calibrada-ch}\\
 \operatorname{Ext}^{\rm fib,cal}_{e_{\varepsilon,n,r}}
   &:=\{(\pi_{\mathcal F}x_{\varepsilon,r-1}(x),
          \pi_{\mathcal F}x_{\varepsilon,r}(x)):x\in G_n\}.
 \label{hmtfg:eq:ext-fibra-calibrada-ch}
\end{align}
L'extension enrichie synchronisée est
\begin{equation}
 \operatorname{Ext}^{\rm enr,cal}_{e_{\varepsilon,n,r}}
 :=\{(x_{\varepsilon,r-1}(x),x_{\varepsilon,r}(x)):
       x\in G_n\}.
 \label{hmtfg:eq:ext-enriquecida-calibrada-ch}
\end{equation}
La troisième application est définie par
\begin{equation}
 \operatorname{graph}
 \bigl(\operatorname{Upd}^{\rm cal}_{e_{\varepsilon,n,r}}\bigr)
 :=\{(z_{\varepsilon,r}(x),x_{\varepsilon,r}(x)):
       x\in G_n\},
 \label{hmtfg:eq:grafo-upd-calibrado-ch}
\end{equation}
où la seconde composante est exactement celle de
\eqref{hmtfg:eq:upd-unaria-tres-vueltas-ch}. Par conséquent,
\begin{equation}
 \mathcal U^{(1),\rm cal}_{e_{\varepsilon,n,r}}
 :=\operatorname{Upd}^{\rm cal}_{e_{\varepsilon,n,r}}\circ
   \mathsf{Tra}^{\rm cal}_{e_{\varepsilon,n,r}}\circ
   \mathsf{Sel}^{\rm cal}_{e_{\varepsilon,n,r}}
 =\{(x_{\varepsilon,r-1}(x),x_{\varepsilon,r}(x)):
      x\in G_n\}.
 \label{hmtfg:eq:transicion-cal-en-tpk-ch}
\end{equation}
Chaque opérateur est indexé par l’arête
\(e_{\varepsilon,n,r}\) produite par la réduction
\(\mathfrak p_\varepsilon\), et non par un trit fourni de l’extérieur. La
coordonnée \(b_{\rm ref}\) distingue en outre leurs codomaines. Par conséquent,
\(\mathsf{Sel}^{\rm cal}_{e}\),
\(\mathsf{Tra}^{\rm cal}\) et \(\operatorname{Upd}^{\rm cal}\) sont des fonctions
sur leurs domaines respectifs, tandis que
l’union
\(\bigcup_{\varepsilon\in\TRIT}
\mathcal U^{(1),\rm cal}_{e_{\varepsilon,n,r}}\) conserve délibérément
les trois sorties de la trisection.

Pour \(\bullet\in\{\APP,{\rm fib},{\rm enr}\}\), les identités et les compositions ne restent pas implicites. On définit
\begin{align}
 \operatorname{Ext}^{\bullet,\rm cal}_{\operatorname{id}_{\widetilde q}}
   &:=\Delta_{\bullet,\widetilde q},\\
 \operatorname{Ext}^{\bullet,\rm cal}_{e_t\cdots e_1}
   &:=\operatorname{Ext}^{\bullet,\rm cal}_{e_t}\circ\cdots\circ
      \operatorname{Ext}^{\bullet,\rm cal}_{e_1}
 \label{hmtfg:eq:ext-palabra-calibrada-ch}
\end{align}
pour tout mot admissible d’arêtes. Les deux clôtures de chemins APP et de
fibre sont formées séparément,
\begin{equation}
 \operatorname{Path}(\operatorname{Ext}^{\APP,\rm cal}),\qquad
 \operatorname{Path}(\operatorname{Ext}^{\rm fib,cal}),
 \label{hmtfg:eq:dos-clausuras-ext-ch}
\end{equation}
et \(\mathsf{Tra}^{\rm cal}\) les synchronise arête par arête.

Enfin, les niveaux d'états complets et leurs troncatures sont fixés par la récurrence elle-même :
\begin{align}
 H^{\rm cal}_{d_{n,r}}
   &:=\{x_{\varepsilon,r-1}(x):
          x\in G_n,\ \varepsilon\in\TRIT\},\\
 H^{\rm cal}_{d_{n,r}+1}
   &:=\{x_{\varepsilon,r}(x):
          x\in G_n,\ \varepsilon\in\TRIT\},\\
 \rho^{\rm cal}_{d_{n,r}+1,d_{n,r}}
   \bigl(x_{\varepsilon,r}(x)\bigr)
   &:=x_{\varepsilon,r-1}(x).
 \label{hmtfg:eq:truncamiento-calibrado-definido-ch}
\end{align}
Le dernier registre conserve \((n,r,\mathfrak p_\varepsilon)\), de sorte que
le prédécesseur du membre de gauche est unique et \(\rho^{\rm cal}\) est bien
définie. Pour \(b>a\), on pose
\(\rho^{\rm cal}_{b,a}:=\rho^{\rm cal}_{a+1,a}\circ\cdots\circ
\rho^{\rm cal}_{b,b-1}\). Cette application retrouve l’état précédent
complet, y compris sa signature ; elle ne prétend pas inverser une coordonnée écrasée
sans consulter le registre.

\begin{proposition}[Réalisation autonome du TPK au moyen de transitions explicites]
\label{hmtfg:prop:modelo-tpk-calibrado-ch}
La famille
\[
\begin{aligned}
 \mathfrak T_{\rm TPK}^{\rm cal}:=\bigl(
 &(\widetilde q_{n,r})_{n,r},\pi_{\rm obs},
   \widetilde F_{\rm obs},M_{\rm ph},\\
 &(\mathcal A_{\widetilde q}^{\rm cal})_{\widetilde q},
   (\mathcal E_{\widetilde q}^{\rm cal})_{\widetilde q},
   (\mathcal F_{\widetilde q}^{\rm cal})_{\widetilde q},
   \operatorname{Ext}^{\APP,\rm cal},
   \operatorname{Ext}^{\rm fib,cal},
   \operatorname{Ext}^{\rm enr,cal},\rho^{\rm cal},\\
 &\mathsf{Sel}^{\rm cal},\mathsf{Tra}^{\rm cal},
   \operatorname{Upd}^{\rm cal},\mathsf C_{\mathsf M},
   \kappa_{\mathsf M},\operatorname{gr}_{\mathsf M},
   H_{\rm Carr},Y_{\rm Carr},\\
 &\operatorname{Sig}^{\rm cal},\pi_{\mathcal F}\bigr).
\end{aligned}
\]
est un modèle enrichi du TPK. En particulier, chaque paire du membre de droite de \eqref{hmtfg:eq:transicion-cal-en-tpk-ch} est une arête TPK effective, et non la conséquence inférée d'une condition seulement nécessaire.
\end{proposition}

\begin{proof}
Les étiquettes \((n,r)\) rendent disjointes les copies des porteurs et fixent
sans ambiguïté chaque origine et destination. La relation
\(\operatorname{Ext}^{\APP,\rm cal}\) transporte la cellule complète et
préserve les deux identités euclidiennes ; la relation
\(\operatorname{Ext}^{\rm fib,cal}\) relie les fibres des états
complets antérieur et postérieur : elle transporte feuillet, orientation, préfixe,
cylindre, frontière et registre et incorpore la retenue, la mémoire et la signature
actualisées. La sélection ne modifie pas le registre et le transport concatène
exactement une entrée et produit la présignature de destination
\eqref{hmtfg:eq:prefirma-transporte-cal-ch}. L’actualisation recalcule la signature au moyen de
l’application typée \(\operatorname{Sig}^{\rm cal}_{\widetilde q}\) à partir du
feuillet, de l’orientation, de la retenue, de la route, de la frontière et de la mémoire déjà déterminés. L’application
\(\operatorname{Upd}^{\rm cal}\) est unaire car l’arête contenue dans
\(z_{\varepsilon,r}(x)\) détermine ses cocycles.

Les identités sont les diagonales et les extensions par mots sont les
compositions de \eqref{hmtfg:eq:ext-palabra-calibrada-ch}. La concaténation des
préfixes et des registres est associative ; la mémoire vérifie la loi de cocycle et
la retenue vérifie
\eqref{hmtfg:eq:ley-carry-correcta-ch}. Par conséquent, les deux parenthésages de trois
arêtes produisent le même état. La formule
\eqref{hmtfg:eq:truncamiento-calibrado-definido-ch} définit la troncature sur
les états complets et l’unicité du dernier registre prouve qu’elle retrouve
l’origine dans tous les champs. Ce sont les lois d’identité, de source,
de but, de composition, de mémoire, de retenue et de troncature du TPK enrichi.
En particulier, les relations d’extension sont déterminées comme graphes
d’applications effectives, et non comme conséquences de simples conditions
nécessaires.
\end{proof}

Chaque \(\gamma_{\varepsilon,k_n}\) comprend exactement neuf applications de \(\operatorname{Upd}^{\rm cal}_{e_{\varepsilon,n,r}}\circ
\mathsf{Tra}^{\rm cal}_{e_{\varepsilon,n,r}}\circ
\mathsf{Sel}^{\rm cal}_{e_{\varepsilon,n,r}}\). La survie sera prouvée au moyen d'une queue construite et ne sera pas utilisée pour définir \(G_n\).

\begin{lemma}[Trisection interne effective d'un tour]
\label{hmtfg:lem:elevacion-catalogal-nonadica}
Pour tout \(n\ge0\), tout \(x\in G_n\) et tout \(\varepsilon\in\{-1,0,+1\}\), l'expression \(\gamma_{\varepsilon,k_n}(x)\) est définie par neuf transitions TPK, avec leurs relations synchronisées \(\operatorname{Ext}^{\APP,\rm cal}\) et \(\operatorname{Ext}^{\rm fib,cal}\), et satisfait
\begin{align}
 \rho^{\rm cal}_{k_n+9,k_n}\gamma_{\varepsilon,k_n}(x)&=x,
 \label{hmtfg:eq:truncamiento-injerto-ch}\\
 g\bigl(\gamma_{\varepsilon,k_n}(x)\bigr)
   &=g(x),&
 \jmath_{k_n+9}\!\left(\operatorname{mem}_{k_n+9}
        (\gamma_{\varepsilon,k_n}(x))\right)
   &=\jmath_{k_n}\!\left(\operatorname{mem}_{k_n}(x)\right)+u,
 \label{hmtfg:eq:fase-memoria-injerto-ch}\\
 c\bigl(\gamma_{\varepsilon,k_n}(x)\bigr)
   &=c(x)+\varepsilon\kappa_{\rm Carr}.&&
 \label{hmtfg:eq:carry-injerto-ch}
\end{align}
\label{hmtfg:eq:tres-vueltas-propiedades-ch}
\label{hmtfg:eq:tres-vueltas-acarreo-ch}
Le registre ordonné des neuf arêtes détermine \(\varepsilon\) de manière unique. Les trois images sont donc disjointes. Chaque extrémité appartient à \(G_{n+1}\) et admet à nouveau les trois applications \(\gamma_{-1,k_{n+1}},\gamma_{0,k_{n+1}},
\gamma_{+1,k_{n+1}}\).
\end{lemma}

\begin{proof}
Les formules \eqref{hmtfg:eq:celula-completa-fase-ch} parcourent une fois les neuf phases et renvoient la copie \(a_{n+1,1}\) de \(a(1)\). À chaque étape, la relation APP transporte le sextuplet complet ; \(\mathsf{Sel}^{\rm cal}\) fixe le type TRIT à partir de la marque de phase \(r\) conservée dans l'état ; et \(\mathsf{Tra}^{\rm cal}\) ajoute exactement un symbole au préfixe, un segment au cylindre et sa paire ordonnée d'extrémités à la frontière. Par récurrence sur \(r\), les tuples explicites \(y_{\varepsilon,r}(x)\) et \(z_{\varepsilon,r}(x)\) satisfont les prédicats des relations indiquées, et la mise à jour unaire produit l'état du niveau suivant. L'identité \(\rho^{\rm cal}_{d_{n,r}+1,d_{n,r}}
(x_{\varepsilon,r}(x))=x_{\varepsilon,r-1}(x)\) est la définition \eqref{hmtfg:eq:truncamiento-calibrado-definido-ch} ; son caractère bien défini découle de l'unicité du dernier registre. La composition des neuf applications de troncature récupère \(x\), ce qui prouve \eqref{hmtfg:eq:truncamiento-injerto-ch} sans inverser séparément les champs écrasés.

La phase \(9\to1\) apparaît une seule fois. Par
\eqref{hmtfg:eq:memoria-arista-ch} et la loi de cocycle, la partie linéaire de mémoire
transporte la mémoire préalable et sa partie translationnelle ajoute exactement \(u\)
dans la limite directe. La compatibilité
\(\jmath_{d+1}\jmath_{d+1,d}=\jmath_d\) prouve
\eqref{hmtfg:eq:fase-memoria-injerto-ch} ; le retour concerne la phase, non
l’état enrichi.

Par \eqref{hmtfg:eq:carry-arista-derivado-ch}, huit arêtes ont un incrément nul et
l’arête qui complète \(\mathfrak p_\varepsilon\) possède l’incrément calculé
\(\chi(\varepsilon,\varepsilon)\kappa_{\rm Carr}\). La loi
\eqref{hmtfg:eq:ley-carry-correcta-ch} et
\eqref{hmtfg:eq:carry-no-parametrico-ch} donnent
\eqref{hmtfg:eq:carry-injerto-ch}. Le registre conserve
\((\mathfrak p_\varepsilon,2\varepsilon,r_3(2\varepsilon),
\chi(\varepsilon,\varepsilon))\) ; les trois valeurs de ce registre sont
distinctes même lorsque l’incrément scalaire est nul. C’est pourquoi les images
ne sont pas identifiées et le décodeur de bloc retrouve un unique
\(\varepsilon\).

Les définitions des arêtes dépendent seulement de la nouvelle phase et concatènent,
au lieu de remplacer, préfixe, cylindre, frontière et registre. Elles sont aussi
invariantes sous
\(\jmath_d\operatorname{mem}_d\mapsto
  \jmath_d\operatorname{mem}_d+nu\)
et sous la translation de la retenue
d’entrée. Par conséquent, la même construction s’applique à l’extrémité de
tout tour. Pour chaque état fini, la suite explicite
\[
 x,\quad\gamma_{0,k_n}(x),\quad
 \gamma_{0,k_{n+1}}\gamma_{0,k_n}(x),\quad\ldots
\]
est une queue infinie compatible. Cela démontre la survie de tous les
états construits et la disponibilité réitérée des trois branches, sans
postuler aucune des deux propriétés. Nous pouvons donc, enfin,
identifier, conformément aux définitions de niveau précédentes,
\begin{equation}
 H_{k_n}^{\rm cal}:=G_n
 \label{hmtfg:eq:subarbol-calibrado-superviviente-ch}
\end{equation}
pour tout \(n\ge0\).
\end{proof}

La clôture précédente distingue les trois horloges impliquées. Neuf émissions
ramènent la phase dans \(C_9\) et font avancer la mémoire. Douze tours, soit
\(108\) émissions, ramènent en outre la position du calendrier observable ;
le registre et la mémoire ne sont pas non plus effacés à ce moment. L’opérateur
\(\mathcal K_{\rm ph}\) agit uniquement ensuite, comme reconnaissance de
mémoire réduite des tours déjà construits ; il n’engendre jamais les arêtes de
\(\operatorname{Ext}^{\rm enr,cal}\).


\subsection{Raffinement des cylindres et lecture ordonnée}
\label{hmtfg:refinamiento-ordenado}

% XI REV11, ch_elevacion_catalogal.tex: bloques de 54 aristas, beta_blk,
% partición efectiva de 729 subcilindros; ch_cierre_cardinal.tex:375--497.
Pour \(\mathbf b=(b_1,\ldots,b_6)\in\FF_3^6\), soit
\(\upsilon:\FF_3\to\{-1,0,+1\}\) le représentant équilibré. Les trois prolongements construits arête par arête définissent
\[
 \Gamma^{[54]}_{\mathbf b,k}
 =\gamma_{\upsilon(b_6),k+45}\circ\cdots\circ
   \gamma_{\upsilon(b_1),k}.
\]
En partant de \(Y_0=\{x_\star\}\), on pose
\[
 Y_{N+1}
 =\{\Gamma^{[54]}_{\mathbf b,1+54N}(x):
          x\in Y_N,\ \mathbf b\in\FF_3^6\}.
\]
La troncature \(\rho^{\rm blk}_{N+1,N}\) supprime les cinquante-quatre dernières arêtes, en récupérant l’état antérieur au moyen du registre de transition.

\begin{proposition}[Codage effectif des blocs de prolongement]
\label{hmtfg:bloques-729}
La lecture des six couples tritiques de chaque bloc définit des bijections
\[
 \beta_N:Y_N\xrightarrow{\ \cong\ }(\FF_3^6)^N
\]
qui commutent avec la troncature. Chaque état de \(Y_N\) possède exactement \(729\) prolongements de bloc dans \(Y_{N+1}\).
\end{proposition}
\begin{proof}
À chaque tour, le registre conserve le couple tritique et sa retenue. Les trois couples produisent les trois sorties distinctes \(\chi(\varepsilon,\varepsilon)=\varepsilon\), et le registre permet de récupérer la branche utilisée, y compris celle de retenue nulle. Les six concaténations produisent donc tous les mots de \(\FF_3^6\). Deux mots différents se distinguent au premier couple enregistré où ils diffèrent. La récupération du préfixe prouve la compatibilité avec la troncature. Par récurrence sur \(N\), chaque mot de longueur \(N\) possède un unique état antécédent ; le mot suivant peut être n’importe laquelle des \(3^6=729\) possibilités.
\end{proof}

Soit \(\Omega_{\Gamma}^{\rm cal}=\varprojlim_NY_N\). Les bijections précédentes induisent
\[
 \Omega_{\Gamma}^{\rm cal}\cong(\FF_3^6)^{\mathbb N}.
\]
La correspondance et son inverse sont continues : chaque coordonnée finie dépend d’un préfixe fini. Les cylindres sont ouverts et fermés. Chaque cylindre de profondeur \(N\) se partitionne en ses \(729\) sous-cylindres de profondeur \(N+1\) ; chacun contient un prolongement infini, obtenu, par exemple, en poursuivant par la branche nulle dans tous les blocs ultérieurs.

% Antecedentes efectivos del continuo conjunto y del universo interno.
% Se imprimen en ambos destinos; la etiqueta de universo no reemplaza su regla.
% Punto de inclusión común: después de APP--TRIT--TPK y los bloques 729,
% antes de la cobertura arquimediana y de la carta interna de Feigenbaum.
% No duplica 00a ni incluye clasificación cardinal posterior.
% La regla semántica transfinita se declara como tal, no como consecuencia
% de la sola existencia de prolongaciones finitas.
\begingroup
\let\HMTFGBaseSection\subsection
\let\HMTFGBaseSubsection\subsubsection
\let\HMTFGBaseSubsubsection\paragraph
\let\section\HMTFGBaseSection
\let\subsection\HMTFGBaseSubsection
\let\subsubsection\HMTFGBaseSubsubsection
% Propietario reproducido por unidad demostrativa; véase PROPIETARIOS_CONTINUO_UNIVERSO.json.
% Origen: XI ES CUMPLIMIENTO_EDITORIAL_20260928, sections/continuo_conjunto.tex:1--629.
% Cuerpos y pruebas conservados; sólo namespace, rutas y remisiones contextuales declaradas.
\section{Construction conjointe du continu au moyen d'histoires et de prolongations compatibles}
\label{hmtfg:base:iv:continuo-conjunto}
L’arithmétique exprimée par une forme normale conserve une lecture précise de l’état qui la produit. Pour situer cette lecture dans HMT, il est nécessaire de maintenir aussi les histoires, les feuillets et les opérations qui agissent sur eux. Ce chapitre réunit ce substrat commun. Ses différentes constructions s’obtiennent à partir du même arbre d’exécutions APP–TRIT–TPK ; les découpages de l’exposition permettent d’étudier leurs propriétés sans les transformer en domaines indépendants choisis par une application ultérieure.

Le point de départ est le noyau formel exposé précédemment. Une émission admissible contient l’opération de APP, l’orientation du TRIT, le changement de feuillet, le résidu et le quotient, ainsi que la mise à jour effective de mémoire du TPK. Lorsque seul un résidu est exprimé, l’histoire d’exécution appartient toujours à l’objet de départ. Deux exécutions ayant une lecture finale identique ne sont pas identifiées pour ce seul motif.

% Propietario reproducido por unidad demostrativa; véase PROPIETARIOS_CONTINUO_UNIVERSO.json.
% Origen: XI ES CUMPLIMIENTO_EDITORIAL_20260928, sections/supervivencia_punto_fijo.tex:1--146.
% Cuerpos y pruebas conservados; sólo namespace, rutas y remisiones contextuales declaradas.
% BEGIN REV09_SUPERVIVENCIA_PUNTO_FIJO
\subsection{Survie par élagage des prolongements}
\label{hmtfg:base:corpus9:xi:supervivencia}

La mémoire rétrospective distingue les histoires enregistrées par les
opérations APP--TRIT--TPK ; la compatibilité prospective organise leurs
prolongements admissibles. Soit \(S\) l’espace des états complets
atteignables par ces opérations. Chaque état conserve les feuillets, les résidus,
les quotients, l’orientation, la retenue, la mémoire et les données de frontière ; lorsque la
transition dépend de la profondeur, celle-ci fait partie de l’état.
Soient \(E\subseteq S\times S\) la relation effective de prolongement et
\(A\subseteq S\) l’ensemble admissible déterminé par la compatibilité de
ces registres. Ainsi, l’admissibilité se décide au moyen des opérations et
des conditions de frontière qui engendrent l’histoire.

Pour \(X\subseteq S\), nous définissons
\begin{equation}
 \operatorname{Pre}_{E}(X)
 =\{s\in S:\exists t\in X,\ (s,t)\in E\},
 \qquad
 F_{\mathrm{surv}}(X)=A\cap\operatorname{Pre}_{E}(X).
 \label{hmtfg:base:corpus9:xi:operador-poda}
\end{equation}
La suite d’élagage s’obtient par
\begin{equation}
 V_0=A,\qquad
 V_{n+1}=F_{\mathrm{surv}}(V_n),\qquad
 V_\infty=\bigcap_{n\geq0}V_n.
 \label{hmtfg:base:corpus9:xi:sucesion-poda}
\end{equation}

\begin{proposition}[Survie et plus grand point fixe]
\label{hmtfg:base:corpus9:xi:mayor-punto-fijo}
L’ensemble \(V_n\) se compose exactement des états de \(A\) qui admettent
un prolongement de \(n\) étapes dont les états appartiennent à \(A\). La
suite \((V_n)\) est décroissante. Si chaque état a un nombre fini
de successeurs admissibles, alors
\[
 V_\infty=F_{\mathrm{surv}}(V_\infty)
          =\nu F_{\mathrm{surv}},
\]
où \(\nu F_{\mathrm{surv}}\) désigne le plus grand point fixe pour l’inclusion.
Ses éléments sont précisément les états qui admettent un prolongement
infini compatible.
\end{proposition}

\begin{proof}
Pour \(n=0\), un prolongement de longueur zéro consiste en l'état lui-même de \(A\). Si la caractérisation vaut pour \(n\), un état appartient à \(V_{n+1}\) exactement lorsqu'il appartient à \(A\) et possède un successeur dans \(V_n\) ; préfixer cette arête produit un prolongement admissible de \(n+1\) étapes. Cette équivalence prouve l'affirmation par récurrence. En outre, \(V_1\subseteq V_0\), et la monotonie de \(F_{\mathrm{surv}}\) propage \(V_{n+1}\subseteq V_n\).

Soit \(s\in V_\infty\), et écrivons
\(E(s)=\{t\in S:(s,t)\in E\}\). Comme \(s\in V_{n+1}\) pour tout \(n\),
les ensembles \(E(s)\cap V_n\) sont non vides, finis et décroissants.
Leur intersection contient un élément \(t\), qui appartient à
\(V_\infty\) et vérifie \((s,t)\in E\). Par conséquent,
\(V_\infty\subseteq F_{\mathrm{surv}}(V_\infty)\).
Réciproquement, la monotonie donne
\(F_{\mathrm{surv}}(V_\infty)\subseteq V_{n+1}\) pour tout \(n\) ;
l’image est également contenue dans \(A=V_0\), de sorte qu’elle appartient à
\(V_\infty\). Cela prouve l’égalité de point fixe.

Si \(W=F_{\mathrm{surv}}(W)\), alors \(W\subseteq A=V_0\).
De \(W\subseteq V_n\), on déduit
\(W=F_{\mathrm{surv}}(W)\subseteq F_{\mathrm{surv}}(V_n)=V_{n+1}\).
La récurrence implique \(W\subseteq V_\infty\), qui prouve la maximalité.
Enfin, l’arbre des prolongements d’un état de \(V_\infty\)
est à ramification finie et possède un niveau non vide à toute profondeur.
Le lemme de König fournit une branche infinie. Dans le sens inverse, les
préfixes de tout prolongement infini situent son état initial
dans tous les \(V_n\).
\end{proof}

\paragraph{Ramification et stabilisation.}
L'hypothèse de finitude possède un contrôle explicite. Considérons une racine avec une branche terminale distincte de chaque longueur entière positive, toutes les branches étant admissibles et disjointes après la racine. La racine admet des prolongements de toute longueur finie et appartient à tous les \(V_n\) ; chacun de ses successeurs appartient à une branche de longueur bornée. La racine ne possède donc ni prolongement infini ni successeur dans \(V_\infty\). Cet arbre a une ramification infinie à la racine.
En revanche, si \(A\) est fini, chaque inclusion stricte \(V_{n+1}\subsetneq V_n\) élimine au moins un état. L'élagage atteint un point fixe après, au plus, \(|A|\) éliminations strictes de couches. Cette borne appartient au domaine fini \(A\).

\paragraph{Unicité de la continuation.}
Pour \(s\in V_\infty\), les successeurs qui conservent la survie forment
\[
 E(s)\cap V_\infty.
\]
La sélection de l’état suivant est unique exactement lorsque
\begin{equation}
 \#\bigl(E(s)\cap V_\infty\bigr)=1.
 \label{hmtfg:base:corpus9:xi:unicidad-sucesor}
\end{equation}
Un prolongement infini est unique lorsque cette condition est maintenue
à chaque état de la route. En effet, chaque successeur survivant possède
un prolongement infini par la proposition précédente ; deux successeurs
distincts produisent deux routes distinctes. Lorsque le successeur est unique à
chaque étape, les préfixes sont déterminés successivement.
La pluralité des successeurs viables exprime une multiplicité de
prolongements dans le même domaine de survie.

\begin{proposition}[Conjugaison des deux directions de prolongement]
\label{hmtfg:base:corpus9:xi:conjugacion-supervivencia}
Soit \(C:S\to S\) une involution telle que \(C(A)=A\) et
\begin{equation}
 (x,y)\in E\quad\Longleftrightarrow\quad(Cy,Cx)\in E.
 \label{hmtfg:base:corpus9:xi:involucion-prolongaciones}
\end{equation}
Notons \(V_n^+\) l’élagage pour \(E\) et \(V_n^-\) l’élagage
pour \(E^{-1}\), tous deux d’ensemble initial \(A\). Alors
\[
 C(V_n^+)=V_n^-\quad(n\geq0),\qquad
 C(V_\infty^+)=V_\infty^-.
\]
\end{proposition}

\begin{proof}
Si \(s_0,\ldots,s_n\) est une route admissible pour \(E\), la suite
\(Cs_0,\ldots,Cs_n\) est une route admissible pour \(E^{-1}\) :
pour chaque \(j\), l’hypothèse convertit
\((s_j,s_{j+1})\in E\) en \((Cs_{j+1},Cs_j)\in E\).
Tous ses états appartiennent à \(A\). La même opération appliquée une autre
fois récupère la route initiale car \(C^2=I\) ; elle établit donc une
bijection entre les prolongements finis des deux directions.
La caractérisation de \(V_n^\pm\) par les routes donne
\(C(V_n^+)=V_n^-\). Comme \(C\) est bijective, elle commute avec
l’intersection de ces ensembles, et l’on obtient l’égalité de
\(V_\infty^\pm\).
\end{proof}

L’involution transporte l’ensemble survivant d’une direction vers celui
de la direction conjuguée. La prolongeabilité infinie dans les deux
directions correspond à \(V_\infty^+\cap V_\infty^-\) ; l’unicité
d’une route conserve le critère de
\eqref{hmtfg:base:corpus9:xi:unicidad-sucesor}. La rétention d’une généalogie
et la restriction de son arbre de continuations sont deux opérations
typées sur les mêmes registres. L’attribution d’une entropie à
ces lectures exige, en outre, la mesure d’histoires correspondante.
% END REV09_SUPERVIVENCIA_PUNTO_FIJO


\subsection{Histoires admissibles et mémoire par composition}
Soit \(H_0:=\{\ast\}\). Le niveau \(H_1\) est l’ensemble fini non vide des
germes APP--TRIT admissibles, et, pour \(k\geq1\), \(H_k\) est l’ensemble des
histoires enrichies de longueur \(k\) qui admettent des prolongements TPK
compatibles à tout horizon fini. Ses éléments se construisent par émissions
successives à partir du domaine initial du noyau. Nous écrivons \(\zeta\varepsilon\)
pour le prolongement de \(\zeta\) par une émission admissible \(\varepsilon\), et
\(\rho_k(\zeta\varepsilon)=\zeta\) pour la suppression de la dernière émission. On
conserve l’étiquette produite par l’émission, même lorsque des étiquettes
différentes expriment un même état scalaire. L’espace des histoires complètes
est défini par
\[
 H_\infty^{\rm enr}:=\varprojlim_k(H_k,\rho_k).
\]
Le sous-système cofinal \(X_N:=H_{6N}\), avec \(X_0=H_0\), coïncide avec l'horloge
de sextuplets utilisée dans la preuve de cardinalité ; par cofinalité, il existe une
identification canonique
\(H_\infty^{\rm enr}\cong\varprojlim_N(X_N,\rho_{6(N+1),6N})\).

Dans les deux feuillets arithmétiques, les données locales satisfont les divisions exactes

\[
\delta_\diamond=r_\diamond+9q_\diamond,
\qquad 0\leq r_\diamond<9,
\qquad \diamond\in\{\Sigma,\Pi\}.
\]
La coordonnée orientée du TRIT admet, de manière analogue, une écriture \(\Delta=o_r+3\kappa\), avec le représentant orienté déclaré par le noyau. La conservation du quotient permet de composer ces représentations sans remplacer une émission par son résidu isolé.

Dans un système de coordonnées de mémoire, la mise à jour d’une émission s’écrit

\[
U_\varepsilon(m)=C_\varepsilon m+\kappa_\varepsilon.
\]
Pour une trajectoire \(\alpha=(\varepsilon_1,\ldots,\varepsilon_k)\), l'opération effective est la composition ordonnée de ces mises à jour. Sa partie linéaire et sa translation sont

\[
C_\alpha=C_{\varepsilon_k}\cdots C_{\varepsilon_1},
\qquad
\kappa_\alpha=
\sum_{j=1}^{k}C_{\varepsilon_k}\cdots
C_{\varepsilon_{j+1}}\kappa_{\varepsilon_j}.
\]
Le produit vide de la dernière somme est l'identité.

\begin{proposition}[Compatibilité affine de la concaténation]
\label{hmtfg:base:prop:continuo-compatibilidad-afin}
Si \(\alpha\) précède \(\beta\), alors

\[
C_{\beta\alpha}=C_\beta C_\alpha,
\qquad
\kappa_{\beta\alpha}=C_\beta\kappa_\alpha+\kappa_\beta,
\qquad
U_{\beta\alpha}=U_\beta U_\alpha.
\]
\end{proposition}
\begin{proof}
Substituer \(U_\alpha(m)\) dans \(U_\beta\) produit \(C_\beta C_\alpha m+C_\beta\kappa_\alpha+\kappa_\beta\). La décomposition en partie linéaire et translation donne les trois identités. L'associativité de la composition démontre que le résultat ne dépend pas d'une subdivision auxiliaire de la trajectoire.
\end{proof}

Cette identité de cocycle est le contenu précis de la mémoire accumulée. La quantité transportée par une seconde trajectoire dépend de la mise à jour du système de coordonnées qu’elle réalise ; ce n’est pas, en général, la somme sans transport de deux registres.

Le retour nonadique conserve une seconde distinction. Dans les coordonnées de phase et de tour, avec \(\bar r\in\{0,\ldots,8\}\), sa représentation est

\[
\sigma_9(w,r)=
\left(w+\left\lfloor\frac{\bar r+1}{9}\right\rfloor,
r+1\pmod9\right),
\qquad
\sigma_9^9(w,r)=(w+1,r).
\]
Pendant neuf étapes, il se produit exactement un franchissement de \(8\) à \(0\). La phase revient et la coordonnée de tour augmente. L'holonomie enrichie retient en outre les mises à jour de mémoire qui accompagnent ce parcours.

\subsection{Raffinement de l'arbre et mesure des histoires}
La mesure initiale est la mesure canonique du recensement : \(\mu(\ast)=1\) et \(\mu(c)=1/|H_1|\) pour chaque germe \(c\in H_1\). Dans cette section, chaque histoire possède un nombre fini et positif de prolongements. La suppression d'émissions définit les applications de raffinement. L'admissibilité du noyau fournit les prolongements ; lorsqu'une résidence utilise une émission neutre, celle-ci donne explicitement une section de \(\rho_k\). L'existence de cette section se rapporte à l'arbre qui inclut cette émission, et non à toute restriction arbitraire de routes.

Soit \(d(\zeta)\geq1\) le nombre d'émissions admissibles en \(\zeta\). L'ordre
APP et le recensement des prolongements déterminent la loi canonique
\(p_{\rm can}(\varepsilon\mid\zeta)=1/d(\zeta)\). Sans paramètres
supplémentaires, cette loi définit

\[
\mu(\zeta\varepsilon)=\frac{\mu(\zeta)}{d(\zeta)}.
\]
Une pondération non uniforme ne peut remplacer \(p_{\rm can}\) que lorsqu'elle a déjà
été produite et ajoutée au registre enrichi. Le constructeur corrélatif
de ce chapitre utilise exclusivement \(p_{\rm can}\) ; sa mesure
est donc une fonction de l'arbre APP--TRIT--TPK, et non d'un choix externe.

\begin{proposition}[Compatibilité cylindrique et espérance conditionnelle]
\label{hmtfg:base:prop:continuo-esperanza-condicional}
À chaque niveau, on a

\[
\sum_{\rho_k(y)=x}\mu(y)=\mu(x).
\]
Dans les espaces \(L^2(H_k,\mu_k)\), les applications

\[
(\mathsf J_kf)(y)=f(\rho_k(y)),
\qquad
(E_kg)(x)=\sum_{\rho_k(y)=x}\frac{\mu(y)}{\mu(x)}g(y)
\]
satisfont \(E_k=\mathsf J_k^*\), \(E_k\mathsf J_k=I\) et \(\|\mathsf J_kf\|=\|f\|\).
\end{proposition}
\begin{proof}
La première identité est la normalisation des probabilités locales. En l'appliquant à \(|f(x)|^2\), on obtient l'isométrie. En regroupant la somme par fibres de \(\rho_k\), on obtient la formule de \(E_k\) ; sur une fonction constante dans chaque fibre, cette formule restitue sa valeur.
\end{proof}

Les masses compatibles définissent une mesure de probabilité sur l'espace des histoires infinies : le cylindre d'une histoire finie reçoit sa masse \(\mu(\zeta)\). La mesure s'étend à partir de l'algèbre des cylindres, et son unicité procède du fait que ceux-ci engendrent la \(\sigma\)-algèbre. Si l'arbre est à ramification finie et que chaque histoire possède une prolongation, la limite inverse de ses niveaux par suppression des émissions est non vide. Dans la topologie cylindrique, chaque cylindre non vide a une mesure positive.

Cette mesure distingue les histoires qui convergent vers une même représentation. La mesure uniforme sur les valeurs terminales serait une lecture différente, car elle éliminerait les multiplicités que la construction vient de conserver.

% Propietario reproducido por unidad demostrativa; véase PROPIETARIOS_CONTINUO_UNIVERSO.json.
% Origen: XI ES CUMPLIMIENTO_EDITORIAL_20260928, sections/ch_memoria_y_naturalidad.tex:1--77.
% Cuerpos y pruebas conservados; sólo namespace, rutas y remisiones contextuales declaradas.
\subsubsection{Lecteurs enrichis et mémoire de fibre}
\label{hmtfg:base:sec:ch-memoria-fibra}

La différence entre histoire et valeur admet un critère opérationnel. À un
niveau fini \(H_k\), soient \(q_k:H_k\to Z_k\) un lecteur et
\(M_k:H_k\to\mathcal M_k\) le registre conservé. Le lecteur enrichi
est \(\widehat q_k=(q_k,M_k)\). La mémoire distingue des états d'une même
fibre du lecteur ; sa suffisance se vérifie dans le domaine considéré.
Cette formulation intègre à l'article lui-même le développement complet de
l'information de fibre utilisée ci-après.

\begin{proposition}[Récupérabilité à un niveau fini]
\label{hmtfg:base:prop:ch-recuperabilidad-fibra}
Il existe une inverse de \(\widehat q_k\) sur son image si et seulement si \(\widehat q_k\) est injectif. Pour une variable aléatoire \(Z\) à valeurs dans \(H_k\) et de distribution \(\mu_k\), on a
\[
 \mathsf H(Z)=\mathsf H(q_k(Z))+\mathsf H(Z\mid q_k(Z)).
\]
Si \(\widehat q_k\) est injectif sur le support de \(Z\), alors
\[
 \mathsf H(Z\mid q_k(Z),M_k(Z))=0.
\]
Ici, \(\mathsf H(Z)=-\sum_z\Pr(Z=z)\log\Pr(Z=z)\), avec
\(0\log0=0\), est l'entropie finie en nats ; l'entropie conditionnelle est
la moyenne de l'entropie des distributions conditionnelles.
\end{proposition}

\begin{proof}
Un inverse récupère un unique antécédent, ce qui exige l’injectivité ;
réciproquement, celle-ci permet d’attribuer à chaque paire exprimée son unique
antécédent. Puisque \(q_k(Z)\) est une fonction déterministe de \(Z\),
regrouper la somme qui définit \(\mathsf H(Z)\) par fibres de \(q_k\) donne la
règle de la chaîne. Si la paire \((q_k(Z),M_k(Z))\) distingue le support,
chaque distribution conditionnelle correspondante est concentrée sur un seul
état et son entropie est nulle.
\end{proof}

Ne pas exprimer une coordonnée et effacer le registre qui permet de la récupérer sont
des opérations différentes. Ce résultat n’identifie pas l’entropie à la cardinalité
du continu et n’introduit pas d’interprétation thermique : son domaine, son lecteur et sa
distribution sont ceux qui viennent d’être fixés.

\subsubsection{Équivariance de la mesure sous les symétries de l'arbre}
\label{hmtfg:base:sec:ch-naturalidad-medida}

\begin{proposition}[Transport conjoint de l'arbre et de sa mesure]
\label{hmtfg:base:prop:ch-naturalidad-medida}
Soit \(g=(g_k)_{k\geq0}\) une collection de bijections \(g_k:H_k\to H_k\)
qui commute avec les troncatures, fixe la racine et transporte bijectivement
les émissions admissibles, en conservant leurs multiplicités. Alors
\[
 d(g_k\zeta)=d(\zeta),\qquad
 (g_k)_*\mu_k=\mu_k,\qquad (g_\infty)_*\mu=\mu.
\]
Plus généralement, si \(\mu_\zeta\) est la loi des continuations d'une
histoire \(\zeta\in H_k\), on a
\[
 (g_\infty)_*\mu_\zeta=\mu_{g_k\zeta}.
\]
\end{proposition}

\begin{proof}
La bijection des émissions conserve leur nombre avec multiplicités. La
distribution uniforme des germes est également invariante. Chaque trajectoire
et son image reçoivent donc le même produit de poids locaux
\(1/d(\zeta)\). On obtient l'égalité des mesures à chaque niveau et dans
chaque cylindre. La cohérence avec les troncatures définit \(g_\infty\),
et l'unicité de la mesure déterminée par les cylindres prouve l'égalité
à la limite. Le même argument, en commençant en \(\zeta\), prouve
l'affirmation sur les continuations.
\end{proof}

On obtient ainsi la naturalité de l'arbre et de la mesure : un changement cohérent
de représentation transporte les poids et n'oblige pas à les choisir à nouveau.
La conservation des émissions et des multiplicités est une hypothèse essentielle de
la proposition, et non une donnée ajoutée après la construction.


\subsection{Transport par feuillets et comptabilité exacte des sorties}
Considérons d'abord l'espace dont la base orthonormée est formée des histoires du niveau \(k\). Le transport qui conserve l'étiquette de chaque prolongement est

\[
\mathcal U_ke_\zeta=
\sum_{\varepsilon\ \mathrm{admisible\ en}\ \zeta}
\sqrt{p_{\rm can}(\varepsilon\mid\zeta)}\,e_{\zeta\varepsilon}.
\]
Les fibres successives d'antécédents distincts sont disjointes. La normalisation de \(p_{\rm can}\) démontre donc \(\mathcal U_k^*\mathcal U_k=I\). La représentation au moyen des masses des histoires est reliée à celle-ci par le changement de base diagonal qui incorpore \(\sqrt{\mu(\zeta)}\).

L'étiquette de feuillet fait partie du registre complet d'une histoire. On écrit
\[
\widehat\zeta=((\zeta_0,s_0),\ldots,(\zeta_k,s_k)),\qquad
s_j\in\{\Sigma,\Pi\},
\]
où \(s_j\) est le feuillet actif à l'étape \(j\). La prolongation ajoute l'émission et le feuillet \(s_{k+1}\) déterminé par le sélecteur opératoire du TPK, en conservant le préfixe intégralement, y compris le feuillet initial. Dans la base de ces histoires relevées,
\[
\mathcal U_ke_{\widehat\zeta}
=\sum_{\varepsilon\ \mathrm{admisible\ en}\ \widehat\zeta}
\sqrt{p_{\rm can}(\varepsilon\mid\widehat\zeta)}\,e_{\widehat\zeta\varepsilon},
\qquad
P_ke_{\widehat\zeta}=\mathbf1_{s_k=\Sigma}e_{\widehat\zeta},\quad
Q_k=I-P_k.
\]
Ainsi, \(P_k\) lit le feuillet actuel, tandis que la distinction entre histoires est maintenue par leur préfixe complet. Deux histoires initiales qui arrivent au même feuillet actif conservent des fibres successeures distinctes. La preuve d’isométrie par collections disjointes de prolongements s’applique littéralement à cette base.
L’orientation tritique est une autre coordonnée du registre : la valeur \(0\) conserve la lecture de seuil dans chaque feuillet et n’ajoute pas de troisième feuillet à cette décomposition. Nous séparons maintenant la conservation de feuillet et l’échange de feuillet :

\[
A_k=P_{k+1}\mathcal U_kP_k+Q_{k+1}\mathcal U_kQ_k,
\qquad
\eta_k=Q_{k+1}\mathcal U_kP_k+P_{k+1}\mathcal U_kQ_k.
\]
\begin{proposition}[Identité de Gram par feuillets]
\label{hmtfg:base:prop:continuo-gram-hojas}
Pour un transport linéaire \(U\), isométrique ou non, la décomposition précédente satisfait

\[
A^*A+\eta^*\eta=PU^*UP+QU^*UQ.
\]
En particulier, pour \(U=\mathcal U_k\),

\[
A_k^*A_k+\eta_k^*\eta_k=I.
\]
\end{proposition}
\begin{proof}
Dans \(A^*A\) et \(\eta^*\eta\), les termes mixtes contenant \(P_{k+1}Q_{k+1}\) s'annulent. En additionnant les termes restants apparaît \(P_kU^*(P_{k+1}+Q_{k+1})UP_k\), ainsi que son analogue avec \(Q_k\). La première identité résulte de la complémentarité des projecteurs ; la seconde utilise en outre \(U^*U=I\).
\end{proof}

La première identité est la formule générale. Remplacer son membre de droite par \(U^*U\) requiert l'annulation des blocs croisés de ce gramien. La conservation isométrique du transport des histoires fournit ici cette condition ; elle ne doit pas être attribuée automatiquement à une compression ultérieure.

Écrivons \(F_{0,0}=I\), \(F_{0,n}=A_{n-1}\cdots A_0\) et \(T_N=F_{0,N}^*F_{0,N}\).

\begin{theorem}[Décomposition finie avec résidu terminal]
\label{hmtfg:base:thm:continuo-descomposicion-terminal}
Pour chaque horizon \(N\),

\[
I=\sum_{n=0}^{N-1}
F_{0,n}^*\eta_n^*\eta_nF_{0,n}+T_N.
\]
La suite \(T_N\) est positive et décroissante, et possède une limite forte positive \(T_\infty\). La somme des sorties converge fortement vers \(I-T_\infty\).
\end{theorem}
\begin{proof}
La proposition~\ref{hmtfg:base:prop:continuo-gram-hojas} donne \(T_n-T_{n+1}=F_{0,n}^*\eta_n^*\eta_nF_{0,n}\). La somme télescopique prouve l'égalité finie et la monotonie. Pour chaque vecteur, les formes quadratiques décroissent vers une limite ; la polarisation définit une forme bornée positive et, par le théorème de représentation des formes bornées, un opérateur \(T_\infty\). La convergence est forte parce que, pour un opérateur positif \(S\leq I\), \(\|Sv\|^2\leq\langle v,Sv\rangle\) ; on l'applique à \(S=T_N-T_\infty\).
\end{proof}

Le résidu terminal est conservé jusqu'à la démonstration de son extinction dans le domaine considéré. Cette distinction empêche qu'un retour de phase, à lui seul, se transforme en une hypothèse d'épuisement de toutes les histoires.

\subsection{Bilan orienté et annulation dans le raffinement}
Une histoire porte son orientation cumulée \(o(h)\) et les décomptes de ses visites aux deux feuillets. Le bilan local correspondant est

\[
b(h)=o(h)\bigl(N_\Sigma(h)-N_\Pi(h)\bigr),
\qquad
b^+=\frac{|b|+b}{2},\quad b^-=\frac{|b|-b}{2}.
\]
Lorsqu’un lecteur exprime la moyenne sur les prolongements, on définit \(b_c=Eb_f\). L’orientation reste dans l’argument de \(b_f\) ; elle n’est pas éliminée avant le calcul de la moyenne.

\begin{proposition}[Défaut exact d'annulation]
\label{hmtfg:base:prop:continuo-defecto-cancelacion}
La quantité

\[
\kappa_{\mathrm{can}}=
\frac{E|b_f|-|Eb_f|}{2}
\]
est non négative et satisfait

\[
E(b_f^+)=(b_c)^++\kappa_{\mathrm{can}},
\qquad
E(b_f^-)=(b_c)^-+\kappa_{\mathrm{can}}.
\]
\end{proposition}
\begin{proof}
L'inégalité triangulaire pondérée donne \(|Eb_f|\leq E|b_f|\). En écrivant les parties positive et négative au moyen de \((|b|\pm b)/2\), on obtient les deux égalités.
\end{proof}

L’information omise lorsqu’on exprime seulement le bilan est l’annulation commune des contributions opposées. Cette quantité est déterminée à partir des mêmes prolongements utilisés par le transport précédent. Le facteur \(1/9\) n’apparaît que lorsque la fibre contient neuf prolongements de poids égaux ; l’espérance conditionnelle est la règle qui reste valable pour une ramification et des pondérations générales.

\subsection{Incidence ternaire des feuillets et complexe rectangulaire}
Les deux feuillets déterminent le module \(V=\mathbb F_3e_\Sigma\oplus\mathbb F_3e_\Pi\). Ses quatre directions projectives sont les droites engendrées par \(e_\Sigma,e_\Pi,e_\Sigma+e_\Pi,e_\Sigma-e_\Pi\). Elles donnent six paires non ordonnées. Les huit vecteurs non nuls sont leurs représentants orientés. Ainsi, les dénombrements \(4,6,8\) appartiennent à des opérations distinctes sur un même module de feuillets.

Sur les six positions de paires, nous définissons

\[
(Ba)_{\{L,L'\}}=a_L+a_{L'},
\qquad
s(q)=\sum_{\{L,L'\}}q_{\{L,L'\}}.
\]
Chaque direction intervient dans trois paires ; d'où \(sB=0\) dans \(\mathbb F_3\). Le lecteur \(W(q)=\mathbf1_{s(q)=0}-1/3\) est invariant sous \(q\mapsto q+Ba\). Sur l'espace ambiant uniforme \(\mathbb F_3^6\), ses moments sont

\[
\mathbb EW=0,\qquad
\mathbb EW^2=\frac29,\qquad
\mathbb EW^3=\frac2{27}.
\]
En effet, \(s\) est surjective et chaque fibre contient \(3^5\) éléments. Le lecteur vaut \(2/3\) avec la probabilité \(1/3\) et \(-1/3\) avec la probabilité \(2/3\). Ces moyennes décrivent l'espace ambiant uniforme ; la distribution induite par un sous-ensemble d'histoires doit être calculée avec ses propres multiplicités, déjà conservées dans \(H_k\).

Pour fixer les ports sans perdre les multiplicités, à la profondeur \(k\), on conserve les lectures étiquetées \(\lambda_\Sigma(\zeta)=(\zeta,\Sigma)\) et \(\lambda_\Pi(\zeta)=(\zeta,\Pi)\). Leurs ensembles sont
\[
\mathscr P_k^\Sigma=\{\lambda_\Sigma(\zeta):\zeta\in H_k\},\qquad
\mathscr P_k^\Pi=\{\lambda_\Pi(\zeta):\zeta\in H_k\}.
\]
Ils sont finis et non vides parce que \(H_k\) l’est. Les valeurs, quotients et résidus
s’expriment sur ces ports ; une égalité de valeurs n’identifie pas leurs
étiquettes. Le produit \(\mathscr P_k^\Sigma\times\mathscr P_k^\Pi\) est le
domaine de la lecture rectangulaire, au sein duquel est enregistré le support des
incidences admissibles.

Ce niveau étant fixé, nous écrivons \(I=\mathscr P_k^\Sigma\), \(J=\mathscr P_k^\Pi\) et \(A_m=\mathbb Z/3^m\mathbb Z\). L'ordre lexicographique APP détermine la branche survivante minimale et, sur celle-ci, les ports de référence \(i_0,j_0\). Aucun choix externe n'est ajouté : changer de coordonnées produit un complexe canoniquement isomorphe. Les applications

\[
\kappa(u,v)_{ij}=u_i-v_j,
\qquad
(\delta w)_{ij}=w_{ij}-w_{ij_0}-w_{i_0j}+w_{i_0j_0}
\]
définissent la suite suivante, valable sur l'anneau fini \(A_m\) :

\[
0\longrightarrow A_m\longrightarrow
A_m^I\oplus A_m^J\mathrel{\mathop{\longrightarrow}^{\kappa}}
A_m^{I\times J}\mathrel{\mathop{\longrightarrow}^{\delta}}
A_m^{(I\setminus\{i_0\})\times(J\setminus\{j_0\})}
\longrightarrow0.
\]
\begin{proposition}[Exactitude de l'incidence rectangulaire]
\label{hmtfg:base:prop:continuo-exactitud-rectangular}
La suite précédente est exacte, avec pour première flèche la diagonale constante.
\end{proposition}
\begin{proof}
L'égalité \(u_i-v_j=0\) pour chaque couple impose que toutes les coordonnées de \(u\) et de \(v\) soient une même constante. La substitution directe donne \(\delta\kappa=0\). Si \(\delta w=0\), prenons \(u_i=w_{ij_0}\) et \(v_j=w_{i_0j_0}-w_{i_0j}\). Alors \(u_i-v_j=w_{ij}\). Enfin, toute matrice du dernier terme se prolonge avec une ligne et une colonne de référence nulles ; son image par \(\delta\) est la matrice initiale. La preuve n'utilise que des sommes et des soustractions ; elle n'exige donc pas de diviser par une unité de l'anneau.
\end{proof}

La réduction de \(A_{m+1}\) à \(A_m\), en maintenant \(k,\mathscr P_k^\Sigma,\mathscr P_k^\Pi,i_0,j_0\) fixes, commute avec les deux applications. Les formules de reconstruction sont les mêmes à chaque profondeur modulaire. C'est la compatibilité arithmétique démontrée dans cette section. Le prolongement des histoires \(k\to k+1\) conserve les étiquettes de préfixe ; ses applications entre ensembles de ports sont une donnée distincte du changement d'anneau \(m\to m+1\).

Ces applications s'obtiennent par la suppression des émissions :
\[
\rho_k^\Sigma(\zeta\varepsilon,\Sigma)=(\zeta,\Sigma),\qquad
\rho_k^\Pi(\zeta\varepsilon,\Pi)=(\zeta,\Pi).
\]
La branche survivante minimale de l'ordre APP qui vient d'être construite possède des prolongements compatibles. Ses deux étiquettes définissent des ports de référence compatibles à tous les niveaux. Le transport des coordonnées est la précomposition :
\[
(L_ku)_{i'}=u_{\rho_k^\Sigma(i')},\quad
(L_kv)_{j'}=v_{\rho_k^\Pi(j')},\quad
(L_kw)_{i'j'}=w_{\rho_k^\Sigma(i'),\rho_k^\Pi(j')}.
\]
Pour le dernier module de la suite, on prolonge d'abord la matrice par zéro sur la ligne et la colonne de référence, on applique cette même formule et on restreint au nouveau complément des ports de référence. Cette règle définit \(L_k\) également lorsqu'un nouveau port a pour antécédent le port marqué.

\begin{proposition}[Naturalité du complexe rectangulaire sous prolongement]
\label{hmtfg:base:iv:rectangular-natural}
Les applications précédentes vérifient
\(L_k\kappa_k=\kappa_{k+1}L_k\) et
\(L_k\delta_k=\delta_{k+1}L_k\).
Elles commutent en outre avec la réduction des coefficients
\(A_{m+1}\to A_m\) et se composent conformément à la suppression des préfixes.
\end{proposition}
\begin{proof}
La première identité s'obtient à partir de \(u_{\rho_k^\Sigma(i')}-v_{\rho_k^\Pi(j')}\). Dans la seconde, chacun des quatre termes définissant \(\delta\) est transporté par la même précomposition. La compatibilité des ports marqués identifie leurs lignes et leurs colonnes ; lorsqu'un antécédent est le port marqué, les quatre termes s'annulent et coïncident avec le prolongement par zéro. Toutes les formules sont des sommes et des différences à coefficients entiers, de sorte qu'elles commutent avec la réduction modulaire. La composition de précompositions est la précomposition avec le préfixe composé, ce qui démontre la dernière affirmation.
\end{proof}

\subsection{Complexe de mémoire et compatibilité des incidences}
La même coordonnée accumulée de mémoire \(c_y\in\mathbb Z\), dans un lieu de résidence \(y\), définit un complexe libre entier. Sa réduction à \(A_m\) utilise le même registre entier avant de le réduire ; les profondeurs successives sont donc compatibles. Ses générateurs sont \(f_y\) en degré \(2\), \(e_y,b_y,r_y\) en degré \(1\), et \(g_y\) en degré \(0\) :

\[
\partial f_y=3c_y e_y+b_y,\qquad
\partial e_y=g_y,\qquad
\partial b_y=-3c_y g_y,\qquad
\partial r_y=0.
\]
L'identité \(\partial^2=0\) se vérifie sur \(f_y\) au moyen de \(3c_yg_y-3c_yg_y=0\), et elle est immédiate sur les autres générateurs.

Si une trajectoire \(\alpha:y\to y'\) met à jour la mémoire, son transport dans le complexe envoie les générateurs homonymes sur leurs nouveaux représentants, sauf

\[
\Psi_\alpha(b_y)=b_{y'}+3(c_{y'}-c_y)e_{y'}.
\]
\begin{proposition}[Naturalité du complexe de mémoire]
\label{hmtfg:base:prop:continuo-naturalidad-memoria}
On a \(\partial\Psi_\alpha=\Psi_\alpha\partial\) et \(\Psi_{\beta\alpha}=\Psi_\beta\Psi_\alpha\). Les applications se prolongent linéairement sur l'anneau de coefficients choisi.
\end{proposition}
\begin{proof}
Sur \(f_y\), l'image de son bord est \(3c_ye_{y'}+b_{y'}+3(c_{y'}-c_y)e_{y'} =3c_{y'}e_{y'}+b_{y'}\). Sur \(b_y\), le bord de son image est \(-3c_{y'}g_{y'}+3(c_{y'}-c_y)g_{y'}=-3c_yg_{y'}\). Les deux autres cas sont immédiats. Dans une concaténation, les différences de mémoire s'additionnent télescopiquement, ce qui prouve la seconde identité.
\end{proof}

La relation entre deux représentants d'une même résidence peut également s'écrire comme un bord explicite. Pour une coordonnée entière \(v\), réduite dans le même anneau le cas échéant, soient

\[
z^-=r_y,\qquad z^+=r_y+3c_yv e_y+v b_y,
\qquad h_*=-vf_y.
\]
Alors \(\partial h_*=z^--z^+\). Sous un transport qui met aussi à jour \(v\), la correction \(K_\alpha=(v_{y'}-v_y)f_{y'}\) satisfait \(K_{\beta\alpha}=K_\beta+\Psi_\beta K_\alpha\). Les symboles \(K_\alpha\) de cette homotopie locale ne désignent pas le sceau dodécaphasique \(K\) : ils appartiennent à un autre domaine et leur fonction est spécifiée par l'identité précédente. Leur fonction peut aussi être vérifiée dans l'égalité

\[
z^+_{y'}-\Psi_\alpha z^+_y=\partial K_\alpha.
\]
Les deux membres valent \((v_{y'}-v_y)(3c_{y'}e_{y'}+b_{y'})\).

\begin{proposition}[Homologie et mémoire du complexe local]
\label{hmtfg:base:prop:continuo-homologia-memoria}
Sur \(\mathbb Z\), ou après réduction à \(A_m\), on a

\[
H_0(\Gamma)=H_2(\Gamma)=0,
\qquad H_1(\Gamma)\cong A_m[r_y]
\]
dans la réalisation modulaire, et \(H_1(\Gamma)\cong\mathbb Z[r_y]\) dans la réalisation entière.
\end{proposition}
\Needspace{9\baselineskip}
\begin{proof}
Définissons l'homotopie de degré \(1\) par \(h(g_y)=e_y\), \(h(b_y)=f_y\), et zéro sur les autres générateurs. Soient \(p\) la projection sur le terme engendré par \(r_y\) et \(\iota\) son inclusion. La vérification sur chaque générateur donne

\[
\partial h+h\partial=I-\iota p.
\]
En particulier, sur \(b_y\), les termes \(3c_ye_y\) et \(-3c_ye_y\) s'annulent. Le complexe se rétracte par homotopie sur le module engendré par \(r_y\), concentré en degré \(1\), ce qui démontre l'affirmation.
\end{proof}

Le transport \(\Psi_\alpha\) est un cisaillement de déterminant \(1\). Ainsi, ce complexe conserve explicitement la manière dont \(c_y\) change au niveau des chaînes, mais son homologie locale ne distingue pas les valeurs de cette coordonnée. La mémoire complète demeure dans les étiquettes et dans les transports de l'état source ; elle ne se restitue pas à partir de \(H_1\) isolément. Une torsion numérique dépendant de la mémoire requiert l'assemblage et la lecture déterminantale spécifiés, outre ce complexe local.

\subsection{Assemblage avec les chemins terminaux}
Pour un horizon fini, les résidences et leurs prolongements forment une catégorie de chemins. À chaque résidence \(\nu\) est attribué le module libre \(W(\nu)=A_m[\operatorname{Path}(\nu,\mathrm{Term})]\), qui conserve chaque continuation terminale comme générateur. Une trajectoire initiale agit dans \(W\) par précomposition. Le complexe \(\Gamma(\nu)\) de la section précédente agit de manière covariante au moyen de \(\Psi\).

L'assemblage utilise les deux transports dans le même module bigradué :

\[
\mathcal B_{q,r}=
\bigoplus_{\nu_0\to\cdots\to\nu_q}
W(\nu_q)\otimes_{A_m}\Gamma_r(\nu_0).
\]
On utilise le complexe normalisé des chemins : les chaînes dégénérées contenant des identités sont omises. L'horizon fini borne alors le nombre de prolongements stricts et le degré horizontal. Le bord horizontal supprime un lieu de résidence intermédiaire en composant les flèches adjacentes. À la première extrémité, il transporte \(\Gamma\) ; à la dernière, il précompose \(W\). Avec les signes alternés, on obtient le bord des chemins \(d_{\mathrm{bar}}\). La différentielle totale est

\[
D=d_{\mathrm{bar}}+(-1)^q\partial_\Gamma.
\]
Les identités de faces annulent par paires les termes de \(d_{\mathrm{bar}}^2\). La naturalité de la proposition~\ref{hmtfg:base:prop:continuo-naturalidad-memoria} implique que \(d_{\mathrm{bar}}\) commute avec \(\partial_\Gamma\) ; le facteur \((-1)^q\) annule les termes croisés du carré. Par conséquent, \(D^2=0\). La compatibilité se démontre sur les générateurs de chemins et n'est pas introduite par une étiquette d'appartenance à un complexe.

L'évaluation terminale n'exige pas non plus de choisir une histoire particulière. Pour chaque continuation \(h\), on applique la composition effective de ses mises à jour. Le résultat conserve la paire formée par \(h\) et sa lecture, avec sa multiplicité. Si \(k<\ell<N\), les continuations se décomposent en une union disjointe de préfixes jusqu'à \(\ell\) et de leurs continuations jusqu'à \(N\). La proposition~\ref{hmtfg:base:prop:continuo-compatibilidad-afin} démontre qu'évaluer directement ou évaluer ces deux parties produit la même lecture. On obtient ainsi le carré de naturalité du terminal.

Les modules de chemins, les incidences, le transport de mémoire et leurs différentielles procèdent des mêmes prolongements. Leur assemblage ne consiste pas à ajouter à un état cinq coordonnées qui seraient conservées par définition : leurs relations sont construites par des sommes, des bords et des compositions effectives, et sont vérifiées sur leurs générateurs.

\subsection{Déterminants, limite et lecture arithmétique}
À chaque emplacement, le complexe local libre \(\Gamma\) admet sa droite déterminante graduée

\[
\operatorname{Det}\Gamma=
\bigotimes_r\left(\bigwedge^{\mathrm{rk}\Gamma_r}\Gamma_r\right)^{(-1)^r}.
\]
L'assemblage total du paragraphe précédent possède, quant à lui, les modules \(\operatorname{Tot}_d\mathcal B=\bigoplus_{q+r=d}\mathcal B_{q,r}\). À horizon fini et sur le même anneau, ils sont libres de rang fini et seul un nombre fini d'entre eux est non nul. Leur droite est
\[
\operatorname{Det}(\operatorname{Tot}\mathcal B)=
\bigotimes_{q,r}
\left(\bigwedge^{\operatorname{rk}\mathcal B_{q,r}}
\mathcal B_{q,r}\right)^{(-1)^{q+r}}.
\]
Cette formule résulte de la décomposition en somme directe à chaque degré. Elle distingue la droite du complexe local et celle de l'assemblage avec tous les chemins ; leurs bases sont ordonnées au moyen des étiquettes conservées. La réduction modulaire commute avec ces puissances extérieures parce que les modules et leurs bases proviennent des mêmes modules libres entiers à horizon fixé.
La droite est définie par les modules et leurs degrés, aussi bien sur \(\mathbb Z\) qu'après réduction à \(A_m\) ; l'exposant négatif désigne le module dual. Une torsion numérique non nulle exige en outre de spécifier le complexe sur un corps, les bases et les identifications d'homologie pertinentes. Pour calculer les facteurs de Smith, on utilise la matrice entière produite avant la réduction, et non un relèvement arbitraire d'une matrice modulaire. Un bloc carré entier de déterminant non nul est inversible sur \(\mathbb Q\) ; sa réduction est inversible sur \(A_m\) seulement lorsque ce déterminant est une unité, c'est-à-dire lorsqu'il n'est pas divisible par \(3\). Le raffinement des mémoires et des incidences permet de formuler et de vérifier la stabilité des facteurs invariants pertinents ; ces conditions ne se déduisent pas de l'existence de la droite déterminante.

Le passage à la limite conserve les systèmes projectifs compatibles de résidus, les histoires et leurs multiplicités. Les carrés prouvés dans les parties précédentes restent commutatifs parce que leurs applications sont données par les mêmes formules à chaque niveau. Pour une coordonnée archimédienne, l’étape supplémentaire consiste à exprimer des intervalles cylindriques non vides, compatibles par inclusion et de diamètre tendant vers zéro. Leurs adhérences compactes emboîtées ont une intersection réduite à un seul point : l’existence s’obtient par compacité et l’unicité par la borne du diamètre. C’est la valeur de la lecture limite. Si les intervalles sont semi-ouverts, le point peut appartenir uniquement à leurs adhérences ; la règle de représentation du bord détermine alors son écriture positionnelle. La valeur ne remplace pas l’histoire qui l’a produite. L’autre représentation conserve les incidences du même état. Les deux destinations partagent leur origine sans pour autant recevoir la même information.

% Procedencia: continuo_conjunto.tex:434--444; remisiones contextuales a
% publicaciones posteriores, no premisas de las pruebas del continuo ni de Feigenbaum.
% Se conserva la descripción positiva sin exigir etiquetas de un anfitrión concreto.
Les régions coinductives, le registre dodécaphasique et l’incidence exceptionnelle s’expriment à partir de ce substrat au moyen de leurs applications respectives. La réalisation hilbertienne conserve ensuite les étiquettes de phase et précise sa propre inclusion unitale ; la dualité et les écrans agissent sur les codomaines qui y sont construits. Les opérations conjointes réunies ici transmettent les histoires, les feuilles, l’orientation et la mémoire à ces représentations. L’identification d’une fonctionnelle particulière à une forme analytique externe exige, à son tour, d’écrire son application sur les générateurs et de vérifier ses produits scalaires : les identités de conservation du présent chapitre ne remplacent pas cette composition.

\subsection{Portée mathématique de la construction conjointe}
Les constructions d'état, d'arbre et de mesure ; les opérations corrélatives ; et
l'assemblage terminal avec naturalité et limite sont développés intégralement
dans cet article. Dans ce chapitre, leurs identités de
bilan, d'incidence et de mémoire ont été formulées et démontrées sur les mêmes
générateurs. Les cinq constructions demeurent ainsi les opérations d'un
unique arbre. Chaque spécialisation conserve visibles ses hypothèses exactes :
isométrie du transport, loi canonique des poids, traitement du résidu
terminal et non-singularité du bloc déterminantal lorsque cela s'applique.

Les équations précédentes prouvent les cocycles de mémoire, la compensation
pondérée, l'exactitude rectangulaire et la compatibilité du complexe de
mémoire à tout horizon fini, ainsi que leur passage compatible à la limite.
Une identification analytique ultérieure doit composer son opérateur sur ces
générateurs et vérifier les propriétés supplémentaires de son codomaine ; elle ne
modifie ni ne remplace la construction conjointe démontrée ici.

\iffalse
% Formulación preliminar conservada fuera del ensamblado; sustituida por la
% versión tipada y verificada que se carga al final de este bloque.
\subsection{Émergence conjointe des cinq constructions}
\label{hmtfg:base:sec:cinco-construcciones-continuo-ch}

Les opérations précédentes sont maintenant réunies sans les transformer en cinq branches. Pour une histoire survivante \(h\in X_k\), un horizon \(N>k\) et une profondeur modulaire \(m\geq1\), soit le registre de base commun
\[
 \mathcal B_{k,N}(h):=
 \bigl(\mathscr T_{k,N}(h),
       (\varepsilon_\alpha,s_\alpha,o_\alpha,
        r_{\Sigma,\alpha},q_{\Sigma,\alpha},
        r_{\Pi,\alpha},q_{\Pi,\alpha},
        \operatorname{Carr}_\alpha,
        \kappa_\alpha,\partial_\alpha,
        \operatorname{Surv}_\alpha)_\alpha;
       \operatorname{src},\operatorname{tgt},\circ,\rho\bigr),
 \tag{C1}
 \label{hmtfg:base:eq:base-comun-registros}
\]
où \(\mathscr T_{k,N}(h)\) est le sous-arbre des prolongements de \(h\)
jusqu'à l'horizon \(N\). Chaque composante du registre est une sortie des
divisions APP, de l'orientation TRIT ou des mises à jour TPK décrites
dans les sections précédentes ; aucune ne procède d'un problème classique choisi
ensuite.

Sur cette même base, on définit les cinq objets structurés
\[
\begin{aligned}
 \mathcal F_{\rm sp}(\mathcal B_{k,N})
   &:=(\ell^2(X_j,\mu_j),\mathcal U_j,A_j,\eta_j,T_j)_{k\leq j\leq N},\\
 \mathcal F_{\rm sol}(\mathcal B_{k,N})
   &:=(b,\kappa_{\rm can},C_\alpha,\kappa_\alpha)_{\alpha\subset\mathscr T_{k,N}},\\
 \mathcal F_{\rm gau}(\mathcal B_{k,N})
   &:=(I_j,J_j,B_j,\kappa_j,\delta_j)_{k\leq j\leq N},\\
 \mathcal F_{\rm coh}^{(m)}(\mathcal B_{k,N})
   &:=(\Gamma_y\otimes A_m,\partial,\Psi_\alpha,H_\bullet)_{y,\alpha\subset\mathscr T_{k,N}},\\
 \mathcal F_{\rm det}^{(m)}(\mathcal B_{k,N})
   &:=(\operatorname{Tot}\mathcal B,\operatorname{Det}(\operatorname{Tot}\mathcal B))_{k,N,m}.
\end{aligned}
 \tag{C2}
\]
Les discriminants \({\rm sp},{\rm sol},{\rm gau},{\rm coh},{\rm det}\) nomment respectivement les lectures spectrale, solénoïdale, de jauge, cohomologique et déterminantale d'un unique objet. Ce ne sont pas des domaines d'entrée. Si \(p_T\) oublie la structure ajoutée par le lecteur \(T\) et renvoie le registre \(\mathcal B_{k,N}(h)\), nous écrivons \(\widetilde{\mathcal F}_T=(\mathcal B_{k,N},\mathcal F_T,p_T)\). La sortie corrélative est le produit fibré
\[
\begin{aligned}
 \mathcal J_{k,N}^{(m)}(h):={}&
 \widetilde{\mathcal F}_{\rm sp}
 \mathop{\times}_{\mathcal B_{k,N}(h)}
 \widetilde{\mathcal F}_{\rm sol}
 \mathop{\times}_{\mathcal B_{k,N}(h)}
 \widetilde{\mathcal F}_{\rm gau}\\
 &\mathop{\times}_{\mathcal B_{k,N}(h)}
 \widetilde{\mathcal F}_{\rm coh}^{(m)}
 \mathop{\times}_{\mathcal B_{k,N}(h)}
 \widetilde{\mathcal F}_{\rm det}^{(m)}.
\end{aligned}
 \tag{C3}
 \label{hmtfg:base:eq:constructor-correlativo}
\]
L'égalité des cinq projections sur \(\mathcal B_{k,N}(h)\) oblige à
conserver le même arbre, les mêmes arêtes et le même registre ; c'est pourquoi
\eqref{hmtfg:base:eq:constructor-correlativo} n'est pas un produit cartésien de cinq
objets assemblés rétrospectivement.

La suppression de la dernière couche de l'arbre induit
\(\operatorname{Res}_{N+1,N}\). Les propositions précédentes donnent, sur les
générateurs, les carrés
\[
\begin{gathered}
 \operatorname{Res}_{N+1,N}\mathcal F_T(\mathcal B_{k,N+1})
 =\mathcal F_T(\mathcal B_{k,N}),
 \qquad T\in\{{\rm sp},{\rm sol},{\rm gau},{\rm coh}\},\\
 \operatorname{Det}(\operatorname{Tot}\mathcal B_{k,N+1})
 \simeq
 \operatorname{Det}(\operatorname{Tot}\mathcal K_{N+1})
 \otimes
 \operatorname{Det}(\operatorname{Tot}\mathcal B_{k,N}),\\
 (\partial\Psi_\alpha=\Psi_\alpha\partial),\qquad
 (L_k\kappa_k=\kappa_{k+1}L_k),\qquad
 (L_k\delta_k=\delta_{k+1}L_k).
\end{gathered}
 \tag{C4}
 \label{hmtfg:base:eq:naturalidad-D}
\]
Le changement \(A_{m+1}\to A_m\) commute également avec \(\partial\), \(\Psi\),
\(\kappa\), \(\delta\) et les droites déterminantes. On obtient donc
un unique prosystème dans \((N,m)\).

Soit \(\mathcal X_\chi\) la limite des histoires enrichies compatibles. Pour \(\xi\in\mathcal X_\chi\), désignons par \(\mathfrak L_N(\xi)\) le préfixe ordonné de ses \(N\) émissions complètes. Les fibres conjointes d'horizon sont
\[
 \mathfrak J_{N,\bullet}^{(m)}
 :=\coprod_{h\in X_0}
 \{(\mathfrak L,J):\mathfrak L\text{ est une branche de longueur }N,
       \ J=\mathcal J_{0,N}^{(m)}(h)\}.
 \tag{C5}
 \label{hmtfg:base:eq:fibras-conjuntas-horizonte}
\]
La marque \(\mathfrak L\) ne retient pas une copie tautologique de \(\xi\) : elle est
formée des émissions et des mises à jour produites sur chaque arête. Le
générateur conjoint est
\[
\begin{aligned}
 \operatorname{Gen}_{\rm cont}:\mathcal X_\chi
   &\longrightarrow\varprojlim_{N,m}\mathfrak J_{N,\bullet}^{(m)},\\
 \xi&\longmapsto
   (\mathfrak L_N(\xi),\mathcal J_{0,N}^{(m)}(h_0))_{N,m},\\
 \mathcal C_{\rm cont}^{\rm disc}
   &:=\operatorname{im}(\operatorname{Gen}_{\rm cont}).
\end{aligned}
 \tag{C6}
 \label{hmtfg:base:eq:continuo-discreto}
\]

\begin{theorem}[Émergence corrélative des cinq constructions]
\label{hmtfg:base:thm:correlacion-cinco}
Pour toute histoire survivante \(h\in X_k\), tout horizon
\(N\geq k+1\) et toute profondeur \(m\geq1\), le produit fibré
\eqref{hmtfg:base:eq:constructor-correlativo} existe et est non vide. Ses cinq
composantes émergent du même registre APP--TRIT--TPK, satisfont les
carrés de naturalité \eqref{hmtfg:base:eq:naturalidad-D}, les changements de base et le
passage à la limite. L'image \(\mathcal C_{\rm cont}^{\rm disc}\) contient cinq
aspects inséparables du même continu, et non cinq continus ni cinq
applications indépendantes.
\end{theorem}

\begin{proof}
La prolongeabilité de l'arbre fournit au moins une arête à chaque sommet intérieur. Le registre de cette arête détermine le transport isométrique et sa décomposition par feuillets, le cocycle orienté, les ports d'incidence, le complexe de mémoire et sa droite déterminante. On construit ainsi successivement \(\mathcal F_{\rm sp}\), \(\mathcal F_{\rm sol}\), \(\mathcal F_{\rm gau}\), \(\mathcal F_{\rm coh}^{(m)}\) et \(\mathcal F_{\rm det}^{(m)}\) sur la même base, ce qui prouve la non-vacuité de \eqref{hmtfg:base:eq:constructor-correlativo}. Aucune des cinq constructions n'apparaît comme coordonnée d'entrée.

La proposition de compatibilité affine prouve la composition de la mémoire ;
l'identité de Gram, l'exactitude rectangulaire et la naturalité du complexe
prouvent les quatre premières lignes de \eqref{hmtfg:base:eq:naturalidad-D}. La
décomposition de la bicomplexité selon la dernière couche prouve l'identité
multiplicative de la droite déterminante. Toutes ces opérations se calculent
sur les mêmes générateurs et commutent avec la réduction modulaire. La fibre,
la relation, le nombre, l'orientation, la retenue, la mémoire, la frontière et
la route demeurent dans \(\mathfrak L_N(\xi)\) ; la multisection, les terminaux
et les lecteurs linéaire et probabiliste demeurent dans
\(\mathcal J_{0,N}^{(m)}(h_0)\). Par compatibilité, la limite produit
\eqref{hmtfg:base:eq:continuo-discreto} sans disjonction.
\end{proof}

La reconnaissance ultérieure de l'unique continu construit au moyen de ses cinq
réalisations spectrale, solénoïdale, de jauge, cohomologique et déterminantale
appartient au langage conventionnel postérieur. Elle ne sélectionne pas ses états et
n'intervient pas dans sa génération.
\fi

% Propietario reproducido por unidad demostrativa; véase PROPIETARIOS_CONTINUO_UNIVERSO.json.
% Origen: XI ES CUMPLIMIENTO_EDITORIAL_20260928, sections/continuo_cinco_construcciones_tipado.tex:1--327.
% Cuerpos y pruebas conservados; sólo namespace, rutas y remisiones contextuales declaradas.
\subsection{Émergence conjointe des cinq constructions}
\label{hmtfg:base:sec:cinco-construcciones-continuo-ch}

Pour \(0\leq k\leq N\), soit
\[
 \rho_{N,k}:=\rho_k\circ\rho_{k+1}\circ\cdots\circ\rho_{N-1},
 \qquad \rho_{k,k}:=\operatorname{id}_{H_k},
\]
et, pour une histoire survivante \(\zeta\in H_k\), définissons
\[
 \operatorname{Desc}_{k,N}(\zeta)
 :=\{x\in H_N:\rho_{N,k}(x)=\zeta\}.
\]
L'arbre fini \(\mathscr T_{k,N}(\zeta)\) a pour sommets
\[
 \coprod_{j=k}^{N}\{y\in H_j:\rho_{j,k}(y)=\zeta\}
\]
et pour arêtes uniquement les prolongements TPK \(a:y\to z\) avec \(z\in\rho_j^{-1}(y)\). Ainsi, ni ses sommets ni ses prolongements ne proviennent d'un arbre auxiliaire.

Pour \(x\in\operatorname{Desc}_{k,N}(\zeta)\), la branche marquée \(\lambda_{\zeta,x}\) est la suite
\[
 \zeta=x_k\longrightarrow x_{k+1}\longrightarrow\cdots
 \longrightarrow x_N=x
\]
déterminée par les troncatures de \(x\). Le registre de base commun est
\[
\begin{aligned}
 \mathscr R_{k,N}(\zeta;x):=
 \bigl(&\mathscr T_{k,N}(\zeta),\lambda_{\zeta,x};\\
 &\bigl(\varepsilon_a,s_a,o_a,
   (\delta_{\diamond,a},r_{\diamond,a},q_{\diamond,a})
       _{\diamond\in\{\Sigma,\Pi\}},
   \operatorname{Carr}_a,C_a,t_a,\operatorname{Fr}_a,
   \operatorname{Ruta}_a,\operatorname{Mem}_a,
   \operatorname{Surv}_a\bigr)_a;\\
 &\operatorname{src},\operatorname{tgt},\circ,
   (\rho_j)_{k\leq j<N}\bigr),
\end{aligned}
\tag{C1}
\label{hmtfg:base:eq:base-comun-registros}
\]
où
\[
 U_a(m_0)=C_am_0+t_a,
 \qquad
 \delta_{\diamond,a}=r_{\diamond,a}+9q_{\diamond,a}.
\]
Le symbole \(m_0\) de cette formule est une coordonnée de mémoire ; ce n'est pas la profondeur modulaire \(m\) utilisée plus bas. Le véritable espace de base est
\[
 \mathfrak R_{k,N}(\zeta)
 :=\{\mathscr R_{k,N}(\zeta;x):
       x\in\operatorname{Desc}_{k,N}(\zeta)\}.
\]
La branche marquée n'est pas une entrée ajoutée : c'est la suite des arêtes
effectivement produites qui se termine en \(x\). Le reste du registre conserve
simultanément les fibres latérales de l'arbre ; la marque ne remplace donc pas
la ramification par une trajectoire isolée.

\Needspace{15\baselineskip}
Posons \(R_m:=\mathbb Z/3^m\mathbb Z\), et soit
\(V_j(\zeta):=\operatorname{Desc}_{k,j}(\zeta)\). Le transport, les
projecteurs de feuillet et le résidu terminal se restreignent au sous-arbre par
\[
 \mathcal U_j^\zeta:=
 \mathcal U_j\!\upharpoonright_{\ell^2(V_j(\zeta),\mu_j)},
 \qquad P_j^\zeta:=P_j\!\upharpoonright_{V_j(\zeta)},
 \qquad Q_j^\zeta:=I-P_j^\zeta,
\]
\[
 A_j^\zeta:=P_{j+1}^\zeta\mathcal U_j^\zeta P_j^\zeta
             +Q_{j+1}^\zeta\mathcal U_j^\zeta Q_j^\zeta,
 \qquad
 \eta_j^\zeta:=Q_{j+1}^\zeta\mathcal U_j^\zeta P_j^\zeta
             +P_{j+1}^\zeta\mathcal U_j^\zeta Q_j^\zeta.
\]
Pour \(k\leq n\leq N\), soient
\[
 F_{k,k}^\zeta:=I,
 \qquad
 F_{k,n}^\zeta:=A_{n-1}^\zeta\cdots A_k^\zeta,
 \qquad
 T_{k,n}^\zeta:=(F_{k,n}^\zeta)^*F_{k,n}^\zeta.
\]
Les ports propres du même sous-arbre sont
\[
 \mathscr P_j^\Sigma(\zeta):=
 \{(y,\Sigma):y\in V_j(\zeta)\},
 \qquad
 \mathscr P_j^\Pi(\zeta):=
 \{(y,\Pi):y\in V_j(\zeta)\};
\]
nous désignons par \(B_j^\zeta,\kappa_j^\zeta,\delta_j^\zeta\) les restrictions de l'incidence et du complexe rectangulaire à ces ports, et par \(L_j^\zeta\) la précomposition induite par la suppression de préfixes.

Pour définir sans abréviation le lecteur de chemins, soit
\[
 W_\ell^\zeta(y):=
 R_m[\operatorname{Path}_{\mathscr T_{k,\ell}(\zeta)}
          (y,\operatorname{Term}_\ell)]
\]
et posons
\[
 \mathscr D_{q,r;k,\ell}^{(m)}[\mathscr R]
 :=\bigoplus_{y_0\to\cdots\to y_q\ \text{dans}\
                  \mathscr T_{k,\ell}(\zeta)}
 W_\ell^\zeta(y_q)\otimes_{R_m}
 \bigl(\Gamma_r(y_0)\otimes_{\mathbb Z}R_m\bigr),
 \qquad
 \operatorname{Tot}_d\mathscr D
 :=\bigoplus_{q+r=d}\mathscr D_{q,r;k,\ell}^{(m)}[\mathscr R].
\]
Les formules déjà démontrées produisent, sur chaque
\(\mathscr R\in\mathfrak R_{k,N}(\zeta)\), les cinq objets
\[
\begin{aligned}
 \mathcal F_{\rm sp}[\mathscr R]
 &:=(\ell^2(V_j(\zeta),\mu_j),
       \mathcal U_j^\zeta,A_j^\zeta,\eta_j^\zeta,
       T_{k,j}^\zeta)_{k\leq j\leq N},\\
 \mathcal F_{\rm sol}[\mathscr R]
 &:=(b,\kappa_{\rm can},C_\alpha,\kappa_\alpha)
      _{\alpha\subset\mathscr T_{k,N}(\zeta)},\\
 \mathcal F_{\rm gau}^{(m)}[\mathscr R]
 &:=(\mathscr P_j^\Sigma(\zeta),\mathscr P_j^\Pi(\zeta),
       B_j^\zeta,\kappa_j^\zeta,\delta_j^\zeta;R_m)_{k\leq j\leq N},\\
 \mathcal F_{\rm coh}^{(m)}[\mathscr R]
 &:=(\Gamma_y\otimes R_m,\partial,\Psi_\alpha,H_\bullet)
      _{y,\alpha\subset\mathscr T_{k,N}(\zeta)},\\
 \mathcal F_{\rm det}^{(m)}[\mathscr R]
 &:=(\operatorname{Det}(\operatorname{Tot}
       \mathscr D_{\bullet,\bullet;k,\ell}^{(m)}[\mathscr R]))
       _{k\leq\ell\leq N}.
\end{aligned}
\tag{C2}
\label{hmtfg:base:eq:cinco-lectores-comunes}
\]
Ici, \(\mathscr D_{\bullet,\bullet;k,\ell}^{(m)}[\mathscr R]\) est le
bicomplexe normalisé
de chemins et de mémoire des sections précédentes, restreint au sous-arbre
d’horizon \(\ell\) et réduit sur \(R_m\). Le dernier lecteur conserve la
collection ordonnée de droites déterminantales jusqu’à l’horizon \(N\) ; de cette
manière, sa troncature est une projection canonique et ne présuppose pas de
factorisation déterminantale inexistante. La loi \(\mu\) est la loi canonique
du recensement, et les ports sont fixés au moyen de la branche lexicographiquement minimale
de l’ordre APP. Par conséquent, aucun des cinq lecteurs ne reçoit de pondération
ou de port extérieurs.

Pour
\(T\in\mathfrak T:=\{{\rm sp},{\rm sol},{\rm gau},{\rm coh},{\rm det}\}\),
définissons l'espace total typé
\[
 \widetilde{\mathcal F}_{T;k,N}^{(m)}(\zeta)
 :=\coprod_{\mathscr R\in\mathfrak R_{k,N}(\zeta)}
   \{\mathcal F_T^{(m)}[\mathscr R]\}_{\mathscr R},
 \qquad
 p_T:\widetilde{\mathcal F}_{T;k,N}^{(m)}(\zeta)
      \longrightarrow\mathfrak R_{k,N}(\zeta),
\]
où \(\mathcal F_{\rm sp}^{(m)}:=\mathcal F_{\rm sp}\) et \(\mathcal F_{\rm sol}^{(m)}:=\mathcal F_{\rm sol}\) ; autrement dit, l'exposant \(m\) est inerte pour ces deux lecteurs. La projection \(p_T\) oublie uniquement la structure du lecteur et conserve son registre source. Le produit fibré correct est
\[
\begin{aligned}
 \mathcal J_{k,N}^{(m)}(\zeta):={}&
 \widetilde{\mathcal F}_{\rm sp;k,N}^{(m)}(\zeta)
 \mathop{\times}_{\mathfrak R_{k,N}(\zeta)}
 \widetilde{\mathcal F}_{\rm sol;k,N}^{(m)}(\zeta)
 \mathop{\times}_{\mathfrak R_{k,N}(\zeta)}
 \widetilde{\mathcal F}_{\rm gau;k,N}^{(m)}(\zeta)\\
 &\mathop{\times}_{\mathfrak R_{k,N}(\zeta)}
 \widetilde{\mathcal F}_{\rm coh;k,N}^{(m)}(\zeta)
 \mathop{\times}_{\mathfrak R_{k,N}(\zeta)}
 \widetilde{\mathcal F}_{\rm det;k,N}^{(m)}(\zeta).
\end{aligned}
\tag{C3}
\label{hmtfg:base:eq:constructor-correlativo}
\]
L'égalité des cinq projections oblige à conserver le même arbre, les
mêmes arêtes, la même branche produite et le même registre. Pour chaque
\(\mathscr R\), les cinq constructions déterminent l'élément
\[
 s_{k,N}^{(m)}(\mathscr R)
 :=(\mathcal F_T^{(m)}[\mathscr R])_{T\in\mathfrak T}
 \in\mathcal J_{k,N}^{(m)}(\zeta).
\]
L'existence de cet élément provient des opérations de \eqref{hmtfg:base:eq:cinco-lectores-comunes}, et non de la notation du produit fibré.

La suppression de la dernière couche induit
\[
 \operatorname{res}^{\mathfrak R}_{N+1,N}
 \bigl(\mathscr R_{k,N+1}(\zeta;x)\bigr)
 :=\mathscr R_{k,N}(\zeta;\rho_N(x)).
\]
La restriction de chaque opérateur aux générateurs retenus définit
\[
 r^m_{N+1,N}:\mathcal J_{k,N+1}^{(m)}(\zeta)
 \longrightarrow\mathcal J_{k,N}^{(m)}(\zeta).
\]
Dans le lecteur déterminantal, \(r^m_{N+1,N}\) supprime uniquement la dernière
coordonnée de la collection de droites de
\eqref{hmtfg:base:eq:cinco-lectores-comunes}. La réduction
\(R_{m+1}\twoheadrightarrow R_m\) définit, par changement de base,
\[
 q^N_{m+1,m}:\mathcal J_{k,N}^{(m+1)}(\zeta)
 \longrightarrow\mathcal J_{k,N}^{(m)}(\zeta).
\]
Sur chaque générateur, on a
\[
\begin{gathered}
 r^m_{N+1,N}q^{N+1}_{m+1,m}
 =q^N_{m+1,m}r^{m+1}_{N+1,N},\\
 \partial\Psi_\alpha=\Psi_\alpha\partial,
 \qquad L_j^\zeta\kappa_j^\zeta
       =\kappa_{j+1}^\zeta L_j^\zeta,
 \qquad L_j^\zeta\delta_j^\zeta
       =\delta_{j+1}^\zeta L_j^\zeta,
\end{gathered}
\tag{C4}
\label{hmtfg:base:eq:naturalidad-D}
\]
et les compositions des applications \(r\) et \(q\) satisfont les identités
fonctorielles. La première égalité se vérifie parce que la suppression de la dernière couche et
la réduction des coefficients agissent sur des coordonnées distinctes. Les trois autres
sont les identités de naturalité prouvées précédemment. On obtient un unique
prosystème dans \((N,m)\), et non cinq limites juxtaposées.

À la racine, on définit
\[
 \mathfrak J_{N,\bullet}^{(m)}:=\mathcal J_{0,N}^{(m)}(\ast).
\tag{C5}
\label{hmtfg:base:eq:fibras-conjuntas-horizonte}
\]
La structure discrète conjointe du continu est la limite de ces objets corrélatifs :
\[
 \boxed{
 \mathcal C_{\rm cont}^{\rm disc}
 :=\varprojlim_{\substack{N\geq1\\m\geq1}}
   \bigl(\mathfrak J_{N,\bullet}^{(m)},
         r^m_{N+1,N},q^N_{m+1,m}\bigr).}
\tag{C6}
\label{hmtfg:base:eq:continuo-discreto}
\]
C'est une définition par le prosystème engendré par l'arbre complet ; ce n'est
ni l'image ni le graphe d'une application qui retiendrait l'état d'entrée comme
première coordonnée.

Les registres possèdent une marque terminale produite. Leurs projections compatibles induisent
\[
 \Pi:\mathcal C_{\rm cont}^{\rm disc}\longrightarrow H_\infty^{\rm enr}.
\]
Réciproquement, une histoire complète détermine, par ses préfixes déjà émis, les sections compatibles \(s_{0,N}^{(m)}\). On obtient ainsi, après avoir défini le continu,
\[
 \operatorname{Gen}_{\rm cont}:H_\infty^{\rm enr}
 \longrightarrow\mathcal C_{\rm cont}^{\rm disc},
 \qquad
 \Pi\circ\operatorname{Gen}_{\rm cont}
 =\operatorname{id}_{H_\infty^{\rm enr}}.
\tag{C7}
\label{hmtfg:base:eq:seccion-historias-continuo}
\]
Cette section démontre la récupérabilité ; elle n'est pas utilisée pour définir
\(\mathcal C_{\rm cont}^{\rm disc}\).

\begin{theorem}[Émergence corrélative des cinq constructions]
\label{hmtfg:base:thm:correlacion-cinco}
Pour toute histoire survivante \(\zeta\in H_k\), tout \(N>k\), tout
\(m\geq1\) et tout \(\mathscr R\in\mathfrak R_{k,N}(\zeta)\), les cinq
constructions de \eqref{hmtfg:base:eq:cinco-lectores-comunes} sont définies sur le
même registre, et \eqref{hmtfg:base:eq:constructor-correlativo} contient
\(s_{k,N}^{(m)}(\mathscr R)\). Leurs troncatures, réductions modulaires et
transports satisfont \eqref{hmtfg:base:eq:naturalidad-D}. De plus,
\(\mathcal C_{\rm cont}^{\rm disc}\neq\varnothing\), l'application \(\Pi\)
est surjective et scindée par \(\operatorname{Gen}_{\rm cont}\). Les
cinq composantes sont des aspects inséparables du même continu, et non cinq
domaines ni cinq applications indépendantes.
\end{theorem}

\begin{proof}
La survie qui définit \(H_k\) et la loi de prolongation TPK donnent
\(\rho_k^{-1}(\zeta)\neq\varnothing\) pour chaque sommet. Comme les fibres sont
finies, le lemme de Kőnig produit \(H_\infty^{\rm enr}\neq\varnothing\).

Un registre \(\mathscr R\) étant fixé, le transport des histoires et des feuillets produit \(\mathcal F_{\rm sp}[\mathscr R]\) ; la composition affine et le défaut exact d'annulation produisent \(\mathcal F_{\rm sol}[\mathscr R]\) ; les incidences et la suite rectangulaire produisent \(\mathcal F_{\rm gau}^{(m)}[\mathscr R]\) ; le complexe de mémoire et son transport produisent \(\mathcal F_{\rm coh}^{(m)}[\mathscr R]\) ; et les droites de l'assemblage normalisé de chemins produisent \(\mathcal F_{\rm det}^{(m)}[\mathscr R]\). Chaque construction agit sur les mêmes arêtes générées, et non sur un sous-domaine supposé ni sur l'espace ambiant abstrait de \(729\) mots.

Les compatibilités affine, de transport, d’incidence et de mémoire prouvent
\eqref{hmtfg:base:eq:naturalidad-D}. Tronquer la collection déterminantale de
\eqref{hmtfg:base:eq:cinco-lectores-comunes} commute par
définition avec \(r\) ; son changement de base commute avec \(q\). Les sections
\(s_{0,N}^{(m)}\) sont compatibles, de sorte qu’elles produisent
\(\operatorname{Gen}_{\rm cont}\). L’identité
\eqref{hmtfg:base:eq:seccion-historias-continuo} prouve à la fois la non-
vacuité de la limite, la surjectivité de \(\Pi\) et la récupérabilité de chaque
histoire sans l’introduire comme coordonnée de définition.

La fibre, la relation, le nombre, l'orientation, la retenue, la mémoire, la
frontière et la route demeurent dans le registre commun ; la multisection, les
terminaux et les lecteurs linéaire et probabiliste demeurent dans les cinq
constructions corrélatives. Le passage à la limite conserve donc la clôture des
dépendances et ne disjoint pas le continu.
\end{proof}

La reconnaissance ultérieure de l'unique continu construit au moyen du langage spectral, solénoïdal, de jauge, cohomologique ou déterminantal appartient à la comparaison conventionnelle ultérieure. Elle ne sélectionne pas ses états et n'intervient pas dans sa génération APP--TRIT--TPK.


\endgroup
% Propietario reproducido por unidad demostrativa; véase PROPIETARIOS_CONTINUO_UNIVERSO.json.
% Origen: XI ES CUMPLIMIENTO_EDITORIAL_20260928, sections/ch_cierre_cardinal.tex:26--403.
% Cuerpos y pruebas conservados; sólo namespace, rutas y remisiones contextuales declaradas.
\subsection{Prolongement d’états et extension sémantique transfinie}

Soit

\[
 \Sigma_{\rm HMT}
 =\bigl(\dot R_{\rm APP},\dot R_{\rm TRIT},\dot R_{\rm TPK},
 \dot R_{U},\dot R_{\Gamma_9},\dot R_{\rm Ext},\dot R_\rho,
 \dot R_{\pi},\dot R_{\varphi},\dot R_e,\dot R_{\alpha},
 \dot R_{\rm ar},\dot R_{\rm exc}\bigr)
\tag{33}\label{hmtfg:base:eq:ch-signatura}
\]

la signature \textbf{relationnelle} dénombrable dont le code effectif est \(h\). Chaque \(\dot R_S\) est le graphe de l'opérateur \(S\) sur les codes où il est défini ; on ne présuppose pas qu'une fonction externe soit totale dans chaque niveau. Les graphes \(\dot R_{\pi},\dot R_{\varphi},\dot R_e,\dot R_{\alpha}\) sont les lecteurs de sortie construits par la génération coinductive et la fermeture dodécaphasique ; ils ne reçoivent pas comme arguments les valeurs conventionnelles des constantes. Leur présence dans \(\Sigma_{\rm HMT}\) assure que la clôture conserve l'état complet : elle intègre simultanément les routes exprimant une valeur et l'incidence qui se prolonge vers la réalisation exceptionnelle.

À chaque horizon, le code de l'état enregistre les coordonnées APP--TRIT--TPK, la route, le feuillet, la retenue, la frontière et la mémoire ; à ses étapes de représentation, il enregistre également \(P_N,\Phi_N,E_N,A_N\), où \(A_N\) est le préfixe de \(\alpha_{\rm HMT}\) produit par le même état. Ces coordonnées ne sélectionnent pas la branche de l'extérieur : elles font partie de l'historique codé. Une numérotation primitive récursive des coordonnées finies d'un état étant fixée, tout historique inverse \(x_\infty=(x_N)_{N\geq1}\in X_\infty^{\rm enr}\) détermine le code
\[
 \operatorname{code}(x_\infty)
 :=\{\langle N,j,c\rangle:(x_N)_j\text{ a pour code }c\}\subseteq\omega.
\]
Les équations de cohérence
\(\rho_{N+1,N}(x_{N+1})=x_N\) permettent de décoder de manière unique
l’histoire à partir de ce réel. Nous écrivons \(m_\infty\) pour le code de
l’histoire réalisée et \(m_\infty(n)\) pour sa fonction caractéristique sous la
numérotation fixée. Le code opératoire \(h\) et ce registre complet sont
réunis, au moyen d’une fonction primitive récursive d’appariement, dans le réel

\[
 u_{h,m}:=
 \{\langle0,\langle n,h(n)\rangle\rangle:n<\omega\}
 \cup
 \{\langle1,\langle n,m_\infty(n)\rangle\rangle:n<\omega\}.
\tag{33a}\label{hmtfg:base:eq:ch-codigo-conjunto}
\]

Les deux projections primitives récupèrent \(h\) et \(m_\infty\) à partir de \(u_{h,m}\).

La première extension est le prolongement métrique et opératoriel des états finis. Pour \(N\ge1\), soit \(\widehat{\mathcal X}^{\rm enr,raw}_{6N}\) l'ensemble des états enrichis bruts de longueur \(6N\), et définissons en termes ensemblistes le sous-système survivant par
\[
\begin{aligned}
 X_N
 &:=\widehat{\mathcal X}^{\rm enr,surv}_{6N}\\
 &:=\Bigl\{x\in\widehat{\mathcal X}^{\rm enr,raw}_{6N}:\\[-.2ex]
 &\qquad\forall M>N\ \exists z\in
   \widehat{\mathcal X}^{\rm enr,raw}_{6M}\;\text{ tel que}\\[-.2ex]
 &\qquad\operatorname{tr}_{6M,6N}(z)=x\ \wedge\
   z\text{ satisfait chaque transition TPK intermédiaire}\Bigr\}.
\end{aligned}
\]
et posons
\[
 \rho_{N+1,N}:=\operatorname{tr}_{6N+6,6N}\!\upharpoonright X_{N+1},
\]
et
\[
 \operatorname{Ext}_N(x):=
 \{y\in X_{N+1}:(x,y)\in
 \operatorname{Ext}_{6N+6,6N}\}.
\]
Le niveau racine est inclus sans exception implicite :
\[
 X_0:=\{\ast\},\qquad
 \rho_{1,0}:X_1\to X_0\text{ est constant},\qquad
 \operatorname{Ext}_0(\ast):=X_1.
\]
L’exposant \(\mathrm{surv}\) désigne
l’arbre élagué des états qui possèdent des descendants TPK compatibles à tout
horizon fini. La loi de prolongement TPK établie dans les sections
précédentes prouve que les germes
émis appartiennent à cet arbre ; la ramification finie et le lemme de Kőnig
transforment la compatibilité à tout horizon en une branche infinie. En
particulier,

\[
 \forall N\ge0\;\forall x\in X_N\quad
 \varnothing\ne\operatorname{Ext}_N(x)
 \subseteq\rho_{N+1,N}^{-1}(x).
\tag{34}\label{hmtfg:base:eq:ch-extension-no-vacia}
\]

Par le même lemme de l’arbre à ramification finie, tout germe survivant
possède un prolongement compatible

\[
 x_\infty=(x_N)_{N\ge1}\in
 X_\infty^{\rm enr}:=\varprojlim_N(X_N,\rho_{N+1,N}).
\tag{35}\label{hmtfg:base:eq:ch-limite-inverso}
\]

La seconde extension agit \textbf{ensuite} sur le code complet de cette
histoire réalisée.
Elle ne s’infère pas de la seule existence de l’opérateur fini
\(\operatorname{Ext}_N\) : elle formalise à l’échelle ensembliste le principe HMT selon
lequel seuls sont objets du domaine intégral ceux qui possèdent une dérivation
généalogique depuis l’état et ses règles. L’extension sémantique transfinie
est :

\[
 \begin{aligned}
  \mathfrak G_0(h,m_\infty)
  &:=\operatorname{TC}\bigl(\{\omega,u_{h,m},h,m_\infty\}\bigr),\\
  \mathfrak G_{\beta+1}(h,m_\infty)
  &:=\operatorname{Def}_{\Sigma_{\rm HMT}}
       \bigl(\mathfrak G_\beta(h,m_\infty)\bigr),\\
  \mathfrak G_\lambda(h,m_\infty)
  &:=\bigcup_{\beta<\lambda}\mathfrak G_\beta(h,m_\infty)
       \quad(\lambda\text{ limite}),\\
  \mathfrak G_{\rm HMT}(h,m_\infty)
  &:=\bigcup_{\beta\in\operatorname{Ord}}
       \mathfrak G_\beta(h,m_\infty).
 \end{aligned}
\tag{36}\label{hmtfg:base:eq:ch-clausura-transfinita}
\]

L'hypothèse du continu et une bijection cible ne sont pas introduites parmi les conditions qui définissent l'opérateur successeur ; l'opérateur est donné prospectivement par

\[
 \operatorname{Def}_{\Sigma_{\rm HMT}}(M)
 :=M\cup
 \Bigl\{
   \{x\in M:\langle M,\in,
       (\dot R_S\cap M^{k_S})_{S\in\Sigma_{\rm HMT}}\rangle
       \models\varphi(x,\bar p)\}:
   \ulcorner\varphi\urcorner\in\omega,
   \ \bar p\in M^{<\omega}
 \Bigr\}.
\tag{37}\label{hmtfg:base:eq:ch-operador-def}
\]

Autrement dit, elle exprime simultanément tous les sous-ensembles définissables par une formule finie du langage relationnel HMT et des paramètres déjà produits par la généalogie. Elle ne consulte ni valeur réelle cible, ni constante conventionnelle, ni cardinal cible, ni modèle extérieur. La formule~\eqref{hmtfg:base:eq:ch-operador-def} définit la clôture transfinie sur les germes, les opérateurs, l'état enrichi, le registre et la double projection construits précédemment.

\begin{lemma}[Codage conjoint du programme HMT et du registre réalisé]
\label{hmtfg:base:thm:ch-6}
La paire formée par
le code opératoire \(h\) et le registre \(m_\infty\) est codée par un unique
réel \(u_{h,m}\subseteq\omega\), et toute formule de \(\Sigma_{\rm HMT}\) se
traduit uniformément en une formule du langage \(\{\in,u_{h,m}\}\).

\end{lemma}
\begin{proof}
\(h\) et \(m_\infty\) sont des suites d'entiers naturels : \(h\) code la signature, les tables finies et les programmes des opérateurs ; \(m_\infty\) code un seul historique réalisé, et non l'espace complet des historiques. La formule~\eqref{hmtfg:base:eq:ch-codigo-conjunto} et ses deux décodeurs prouvent la première affirmation. Chaque symbole de~\eqref{hmtfg:base:eq:ch-signatura} désigne sa \textbf{relation de graphe codée dans \(h\)}, jamais le graphe externe de toutes ses évaluations possibles. APP, TRIT, \(U_t\), \(\Gamma_9\), les restrictions \(\rho\) et les extensions finies ont des graphes primitifs récursifs sur des codes entiers et, par conséquent, des formules \(\Delta_0\) absolues entre niveaux transitifs. Les deux représentations s'expriment comme relations de graphe au moyen des formules

\[
 \begin{aligned}
 \dot R_{\rm ar}(x,y)
 &\;\Longleftrightarrow\;
 \forall k\in\omega\;\exists N\;\forall n\ge N\quad
 |\mu_n(x)-y|<2^{-k},\\
 \dot R_{\rm exc}(x,c,z)
 &\;\Longleftrightarrow\;
 \forall N\in\omega\quad z\!\upharpoonright N
   =E_N(x\!\upharpoonright N,c\!\upharpoonright N),
 \end{aligned}
\tag{37c$_0$}\label{hmtfg:base:eq:ch-grafos-publicacion}
\]

où \(\mu_n\) et \(E_N\) sont les fonctions rationnelles/finies dont les codes
figurent dans \(h\). On remplace par récurrence chaque formule atomique contenant
un symbole HMT par sa formule de graphe de paramètre \(u_{h,m}\). Les
connecteurs et les quantificateurs sont conservés. Cette récurrence sur la complexité
syntaxique produit une traduction \(\varphi\mapsto\varphi^*\) telle que, pour
tout niveau transitif \(M\) de~\eqref{hmtfg:base:eq:ch-clausura-transfinita} contenant les paramètres,

\[
 \langle M,\in,(\dot R_S\cap M^{k_S})_S\rangle
   \models\varphi(\bar a)
 \quad\Longleftrightarrow\quad
 M\models\varphi^*(\bar a;u_{h,m}).
\tag{37c}\label{hmtfg:base:eq:ch-eliminacion-uniforme}
\]

Ainsi, \eqref{hmtfg:base:eq:ch-operador-def} ne peut consulter ni oracle, ni table décimale cible, ni relation non codée introduisant une infinité non dénombrable d'objets en une seule étape.
\end{proof}

La sémantique HMT est maintenant définie \textbf{indépendamment de son identification à \(\mathfrak G\)}. Soit \(\mathcal T_0^{\rm HMT}\) l'ensemble des noms canoniques des éléments de \(\mathfrak G_0(h,m_\infty)=
\operatorname{TC}(\{\omega,u_{h,m},h,m_\infty\})\). En particulier, le niveau initial contient le registre réalisé \(m_\infty\), et non toutes les histoires de \(X_\infty^{\rm enr}\) comme paramètres primitifs. Les termes admissibles sont les arbres bien fondés obtenus par

\[
 \begin{aligned}
 \mathcal T^{\rm HMT}_{\beta+1}
 &:={\mathcal T}^{\rm HMT}_\beta\cup
 \{\langle\beta,\ulcorner\varphi\urcorner,
       t_0,\ldots,t_{k-1}\rangle:
       t_i\in\mathcal T^{\rm HMT}_\beta\},\\
 \mathcal T^{\rm HMT}_\lambda&:=
       \bigcup_{\beta<\lambda}\mathcal T^{\rm HMT}_\beta,
 \qquad
 \mathcal T^{\rm HMT}:=
       \bigcup_{\beta\in\operatorname{Ord}}\mathcal T^{\rm HMT}_\beta.
 \end{aligned}
\tag{37a}\label{hmtfg:base:eq:ch-terminos}
\]

La valeur d’un terme successeur est le sous-ensemble du niveau précédent défini
par \(\varphi\) dans la structure relationnelle induite, avec les valeurs de
\(t_0,\ldots,t_{k-1}\) comme paramètres. On définit, avant toute
comparaison ensembliste,

\[
 \operatorname{Ob}_{\rm sem}^{\rm HMT}(h,m_\infty)
 :=\{\llbracket t\rrbracket:t\in\mathcal T^{\rm HMT}\}.
\tag{37b}\label{hmtfg:base:eq:ch-objetos-semanticos}
\]

C'est la règle d'\textbf{exhaustivité générative} adoptée par HMT : tout objet du domaine intégral doit posséder l'un de ces arbres de dérivation, et tout arbre admissible exprime un objet. Elle n'est pas une conséquence de la ramification finie. L'hypothèse du continu n'est pas introduite dans cette règle ; elle est obtenue par la démonstration développée dans l'article. Ni une énumération des réels ni une bijection cible ne sont introduites comme données.

\begin{theorem}[Équivalence entre la sémantique des termes et la hiérarchie généalogique]
\label{hmtfg:base:thm:ch-6a}
L’évaluation des
termes précédents vérifie

\[
 \operatorname{Ob}_{\rm sem}^{\rm HMT}(h,m_\infty)
 =\mathfrak G_{\rm HMT}(h,m_\infty).
\tag{37d}\label{hmtfg:base:eq:ch-equivalencia-semantica}
\]

\end{theorem}
\begin{proof}
Par récurrence transfinie. Au niveau initial, les deux classes coïncident. Si toute valeur d'un terme de \(\mathcal T^{\rm HMT}_\beta\) appartient à \(\mathfrak G_\beta\), la sémantique du terme successeur est l'un des sous-ensembles de~\eqref{hmtfg:base:eq:ch-operador-def}. Réciproquement, chaque sous-ensemble de~\eqref{hmtfg:base:eq:ch-operador-def} est donné par un code de formule et un tuple fini de paramètres ; d'après l'hypothèse de récurrence, ces paramètres possèdent des termes et~\eqref{hmtfg:base:eq:ch-terminos} forme le terme qui le désigne. À une limite, les deux constructions prennent la même réunion.
\end{proof}

\Needspace{9\baselineskip}
L’encodeur d’états est compatible avec les troncatures et envoie chaque
prolongement TPK sur un terme successeur. Soit
\(\mathcal T_N^{\rm st}\subset\mathcal T^{\rm HMT}\) la sous-classe de termes
qui code des états complets de profondeur \(N\), et soit
\(\tau_{N+1,N}:\mathcal T_{N+1}^{\rm st}\to\mathcal T_N^{\rm st}\) l’effacement
du dernier bloc avec la troncature correspondante du registre généalogique,
de la retenue et de la mémoire. Si \(J_N:X_N\to\mathcal T_N^{\rm st}\) est l’encodeur
complet, la compatibilité de \(\operatorname{Ext}_N\) prouve l’identité
correctement typée

\[
 \forall x\in X_N\ \forall y\in\operatorname{Ext}_N(x)\qquad
 \tau_{N+1,N}\bigl(J_{N+1}(y)\bigr)=J_N(x).
\tag{37e$_0$}\label{hmtfg:base:eq:ch-naturalidad-codificador}
\]

Définissons au sein de la sémantique
\[
 \operatorname{Succ}^{\rm TPK}_{N+1}(J_N(x)):=
 \{t\in\mathcal T_{N+1}^{\rm st}:\exists y\in X_{N+1}\,
   (y\in\operatorname{Ext}_N(x)\ \wedge\ t=J_{N+1}(y))\}.
\]
Choisissons \(\beta\) tel que \(\mathfrak G_\beta\) contienne \(x\), \(X_N\),
\(X_{N+1}\), les graphes finis de
\(\operatorname{Ext}_N,J_N,J_{N+1}\) et leurs codes dans \(h\). Alors la
formule primitive récursive du graphe de \(\operatorname{Ext}_N\) donne

\[
 J_{N+1}[\operatorname{Ext}_N(x)]
 =\operatorname{Succ}^{\rm TPK}_{N+1}(J_N(x)),
 \qquad
 J_{N+1}[\operatorname{Ext}_N(x)]
 \in\mathfrak G_{\beta+1},
\tag{37e}\label{hmtfg:base:eq:ch-sucesores-semanticos}
\]

et les deux membres sont formés par la formule de graphe de la même extension
TPK. Les équations~\eqref{hmtfg:base:eq:ch-naturalidad-codificador}--\eqref{hmtfg:base:eq:ch-sucesores-semanticos} sont le lien sémantique explicite : la
sous-classe des termes successeurs est exactement l’image des prolongements
admissibles et leur troncature commute. On n’identifie pas la fibre complète de
\(\rho\) à \(\operatorname{Ext}_N\), et l’on n’identifie pas
\(\operatorname{Def}\) à TPK par nomenclature.

\begin{lemma}[Transitivité, axiomes et bon ordre du modèle généalogique]
\label{hmtfg:base:thm:ch-6b}
\(\mathfrak G_{\rm HMT}(h,m_\infty)\) est une classe transitive, contient tous les ordinaux, satisfait la théorie des ensembles de Zermelo--Fraenkel (ZF) et possède un bon ordre global définissable ; elle satisfait donc le choix et tous les axiomes utilisés dans la comparaison cardinale ultérieure.

\end{lemma}
\begin{proof}
Abrégeons \(B=\mathfrak G_0\) et \(u=u_{h,m}\). La récurrence sur~\eqref{hmtfg:base:eq:ch-clausura-transfinita} prouve simultanément que chaque niveau est transitif et que \(\mathfrak G_\alpha\in\mathfrak G_{\alpha+1}\) ; si tous les ordinaux inférieurs à \(\alpha\) sont présents, leur ensemble est définissable au niveau correspondant et exprime \(\alpha\). L'union contient donc tous les ordinaux. L'extensionnalité et la fondation sont héritées de l'univers ambiant ; l'ensemble vide et l'infini appartiennent à \(B\). Un niveau contenant les paramètres étant choisi, les formules définissant \(\{a,b\}\), \(\bigcup a\) et les sous-ensembles par séparation les expriment à un niveau successeur.

Pour la collection--remplacement, soit \(F\) une relation fonctionnelle définissable dans \(\mathfrak G\) sur l'ensemble \(a\). Le schéma de réflexion relatif à la hiérarchie définissable \((\mathfrak G_\xi,\in,u_{h,m})_{\xi\in\mathrm{Ord}}\), appliqué dans la métathéorie ambiante à la liste finie de formules qui définit \(F\), produit un niveau \(\mathfrak G_\theta\) qui contient \(a\), les paramètres, et qui est correct pour ces formules. Les rangs des témoins de \(F[a]\) sont alors bornés par \(\theta\) ; la séparation sur \(\mathfrak G_\theta\) forme l'image dans \(\mathfrak G_{\theta+1}\). Pour l'ensemble des parties, on ne présuppose pas l'ensemble interne que l'on veut construire. Une fois \(a\in\mathfrak G_\gamma\) fixé, on considère dans la métathéorie ambiante l'ensemble déjà existant \(\mathcal P^V(a)\). Pour chaque \(b\in\mathcal P^V(a)\cap\mathfrak G\), soit \(r(b)\) sa première étape généalogique. Comme \(\mathfrak G\) est une classe définissable, la séparation dans la métathéorie ambiante forme d'abord l'ensemble \(\mathcal P^V(a)\cap\mathfrak G\). Le remplacement \textbf{dans la métathéorie ambiante} appliqué à cet ensemble borne les ordinaux \(r(b)\) par un certain \(\theta\). Par conséquent,
\[
 \mathcal P^{\mathfrak G}(a)
 =\{b\in\mathfrak G_\theta:b\subseteq a\},
\]
et le membre droit s'exprime par séparation dans \(\mathfrak G_{\theta+1}\). Il n'introduit pas dans \(\mathfrak G\) les sous-ensembles extérieurs à la sémantique~\eqref{hmtfg:base:eq:ch-objetos-semanticos}. C'est l'argument non circulaire de majoration par les rangs.

Enfin, on attribue à chaque objet son premier niveau de naissance et son
plus petit \textbf{nom de dérivation} : un arbre de dérivation bien fondé et
à ramifications finies (et, par~\eqref{hmtfg:base:eq:ch-terminos}, fini pour chaque nom individuel), étiqueté par un
ordinal de niveau, un code de formule et le tuple des noms de ses paramètres.
Par récurrence transfinie, on ordonne ces noms : d’abord l’étiquette ordinale,
puis le code de formule, puis, lexicographiquement, les noms des
paramètres déjà ordonnés. Le plus petit nom de chaque valeur définit une relation de
classe \(<_{\mathfrak G}\) qui compare tous les objets. C’est un bon ordre
global définissable qui, en prenant le plus petit élément de chaque ensemble non vide,
prouve le Choix.
\end{proof}

% Propietario reproducido por unidad demostrativa; véase PROPIETARIOS_CONTINUO_UNIVERSO.json.
% Origen: XI ES CUMPLIMIENTO_EDITORIAL_20260928, sections/ch_reflexion_testigos.tex:1--89.
% Cuerpos y pruebas conservados; sólo namespace, rutas y remisiones contextuales declaradas.
\paragraph{Réflexion par clôture des témoins et construction de l'image.}
L’étape de Collection--Remplacement du lemme~\ref{hmtfg:base:thm:ch-6b} admet la
démonstration explicite suivante. On maintient la hiérarchie déjà engendrée à partir du
code opératoire et du registre réalisé ; nous abrégeons
\(\mathfrak G=\mathfrak G_{\rm HMT}(h,m_\infty)\).
L’argumentation utilise la Séparation et le Remplacement dans la métathéorie ambiante,
sans présupposer ces schémas dans \(\mathfrak G\).

Soit \(\Phi\) une collection finie de formules. On traduit d'abord leurs quantificateurs universels au moyen de la négation et de la quantification existentielle, puis on ferme la collection résultante par sous-formules. Pour tout niveau initial, il existe un ordinal limite ultérieur \(\theta\) tel que
\[
 \mathfrak G_\theta\models\varphi(\bar a)
 \quad\Longleftrightarrow\quad
 \mathfrak G\models\varphi(\bar a)
 \qquad
 (\varphi\in\Phi,\ \bar a\in\mathfrak G_\theta^{<\omega}).
\]
La notation de satisfaction de la classe est interprétée, pour chaque formule fixée, en relativisant ses quantificateurs à la classe définissable \(\mathfrak G\) ; aucun prédicat uniforme de vérité n'est introduit.

\begin{proof}
Fixons \(\alpha\). Pour chaque sous-formule existentielle
\(\exists y\,\psi(y,\bar x)\) de \(\Phi\), les tuples
\(\bar a\in\mathfrak G_\alpha^{<\omega}\) de l’arité correspondante
qui possèdent un témoin dans \(\mathfrak G\) forment un ensemble par Séparation
métathéorique. À chacun est attribué le plus petit ordinal \(\beta\) pour lequel
un \(y\in\mathfrak G_\beta\) satisfait
\(\mathfrak G\models\psi(y,\bar a)\). L’ordinal existe car
\(\mathfrak G\) est l’union de ses niveaux. Le Remplacement dans la métathéorie
borne ces premières étapes des témoins. La finitude de \(\Phi\) permet
de choisir de manière définissable une seule borne \(f(\alpha)>\alpha\), valable pour
toutes ses sous-formules existentielles. Par conséquent, chaque tuple considéré possède
\emph{au moins un témoin} dans \(\mathfrak G_{f(\alpha)}\).

Nous prenons \(\alpha_0\) suffisamment grand pour contenir les paramètres
initiaux et posons
\[
 \alpha_{n+1}=f(\alpha_n),\qquad
 \theta=\sup_{n<\omega}\alpha_n.
\]
Par continuité de la hiérarchie, tout tuple fini de
\(\mathfrak G_\theta\) appartient à un certain \(\mathfrak G_{\alpha_n}\).
Si une formule existentielle de \(\Phi\) avec ce tuple est vraie dans
\(\mathfrak G\), elle possède un témoin dans
\(\mathfrak G_{\alpha_{n+1}}\subseteq\mathfrak G_\theta\).
La récurrence sur les formules prouve maintenant l’équivalence : les formules atomiques
sont conservées dans la structure induite ; les connecteurs préservent
l’équivalence ; l’étape existentielle utilise ce témoin dans un sens et le
même témoin, vu dans \(\mathfrak G\), dans l’autre. La traduction préalable
des universelles inclut aussi les formules existentielles qui expriment
leurs contre-exemples possibles. La réflexion finie requise est prouvée.
\end{proof}

Soit maintenant \(F(x,y,\bar p)\) fonctionnelle sur \(a\) dans \(\mathfrak G\). Nous appliquons la construction aux formules qui expriment \(F\), l'existence de ses témoins et son image, avec \(a,\bar p\in\mathfrak G_\theta\). La transitivité donne \(a\subseteq\mathfrak G_\theta\). Chaque \(x\in a\) possède sa valeur à l'intérieur de ce niveau, et la réflexion identifie l'image exacte à
\[
 b=\{y\in\mathfrak G_\theta:
       \mathfrak G_\theta\models
       \exists x\in a\ F(x,y,\bar p)\}.
\]
Cet ensemble est définissable sur \(\mathfrak G_\theta\), avec des paramètres du même niveau. Par l'opération successeur de la hiérarchie, \(b\in\operatorname{Def}(\mathfrak G_\theta)
\subseteq\mathfrak G_{\theta+1}\). On obtient ainsi le remplacement dans \(\mathfrak G\). Si l'on conserve l'existence de témoins et que l'on omet l'unicité, la même construction fournit un ensemble de témoins pour la collection.

\paragraph{Arité des noms et sélection interne.}
Dans le bon ordre du lemme~\ref{hmtfg:base:thm:ch-6b}, l’énumération de la base
attribue à chaque élément son plus petit code naturel. À chaque successeur, on ordonne
d’abord les codes de formule ; une fois l’un fixé, l’arité de son
tuple de paramètres est fixée. L’ordre lexicographique s’applique alors à un
\emph{produit fini d’arité fixe} des bons ordres précédents.
Le plus petit nom représente chaque nouvelle valeur et les niveaux limites conservent
l’ordre de première apparition. L’ordre des tuples est ainsi spécifié,
sans mélanger des arités différentes dans un unique ordre
lexicographique. La sélection du plus petit élément de chaque ensemble non vide,
avec le Remplacement déjà obtenu, forme intérieurement les fonctions de
choix utilisées ensuite. Cette construction de noms précède les
injections cardinales et ne reçoit aucune bijection du continu comme
donnée.




\subsection{Lecture ordonnée des cylindres de bloc}
\label{hmtfg:publicacion-cilindros-bloque}

Pour le représentant \([b]_3\in\{0,1,2\}\), on fixe
\[
 \nu(\mathbf b)=\sum_{i=1}^6[b_i]_3\,3^{6-i}\in\{0,\ldots,728\}.
\]
Dans la complétion ordonnée HMT, un préfixe
\(s=(m;\mathbf b_0,\ldots,\mathbf b_{N-1})\), \(m\in\ZZ\), détermine
\[
 a_N(s)=m+\sum_{j=0}^{N-1}\nu(\mathbf b_j)\,729^{-j-1},\qquad
 I_N(s)=[a_N(s),a_N(s)+729^{-N})_{\rm HMT}.
\]
La retenue conserve le raccord entre les deux représentants d’une frontière. Si \(C_N(s)=\overline{I_N(s)}\), les relations précédentes donnent
\[
 I_N(s)=\bigsqcup_{\mathbf b\in\FF_3^6}I_{N+1}(s;\mathbf b),
 \qquad \operatorname{diam}C_N(s)=729^{-N}\longrightarrow0.
\]

\begin{proposition}[Couverture de la lecture archimédienne]
\label{hmtfg:cobertura-arquimediana}
L'application
\[
 P_{\rm ar}:\ZZ\times\Omega_\Gamma^{\rm cal}
       \longrightarrow\mathbb R_{\rm HMT},\qquad
 \{P_{\rm ar}(m,\xi)\}
    =\bigcap_N C_N(m;\beta_N(\rho_{\infty,N}\xi))
\]
est bien définie, est surjective et détermine
\[
 (\ZZ\times\Omega_\Gamma^{\rm cal})/\!\sim_{P_{\rm ar}}
       \ \cong\ \mathbb R_{\rm HMT}.
\]
\end{proposition}
\begin{proof}
Les extrémités gauches forment une suite de Cauchy dans la complétion HMT, puisque les intervalles sont emboîtés et que leurs diamètres tendent vers zéro. Leur limite appartient à tous les intervalles fermés et est unique. Pour tout élément de la complétion, on choisit d’abord l’intervalle entier qui le contient, puis l’unique sous-intervalle semi-ouvert de chaque partition qui le contient. La proposition précédente réalise chacun de ces choix par un prolongement TPK effectif. Les préfixes compatibles définissent \(\xi\), dont la lecture est l’élément choisi. Aux frontières, la retenue identifie les écritures équivalentes après leur construction. Le quotient par égalité de lecture est, par définition et surjectivité, en bijection avec la complétion.
\end{proof}

La construction est également effectuée dans l’univers généalogique interne \(\mathfrak G=\mathfrak G_{\rm HMT}(h,m_\infty)\). Ici, \(h\) code les règles operatorielles et \(m_\infty\) l’histoire compatible ; l’univers s’obtient par clôture transitive initiale, définissabilité à chaque successeur et réunion aux stades limites :
\[
 G_0=\operatorname{TC}\{\omega,u_{h,m},h,m_\infty\},\qquad
 G_{\beta+1}=\operatorname{Def}_{\Sigma_{\rm HMT}}(G_\beta),\qquad
 G_\lambda=\bigcup_{\beta<\lambda}G_\beta .
\]
Les relations de \(\Sigma_{\rm HMT}\) sont celles des opérations déjà construites et de leurs codes. Le domaine interne est donc maintenu dans toutes les affirmations de couverture. Notons \(\mathbb R_{\rm HMT}^{\mathfrak G}\) sa complétion et \(\mathbb R^{\mathfrak G}\) la réalisation ordonnée au sein du même univers.

\begin{proposition}[Carte ordonnée interne]
\label{hmtfg:carta-interna}
Il existe un unique isomorphisme de corps ordonnés archimédiens complets
\[
 \iota:\mathbb R_{\rm HMT}^{\mathfrak G}
             \xrightarrow{\ \cong\ }\mathbb R^{\mathfrak G}
\]
qui fixe l’unité. Il conserve les extrémités rationnelles, les racines carrées positives et les limites de suites convergentes.
\end{proposition}
\begin{proof}
L’unité fixe les entiers et les rationnels. À chaque élément de la complétion HMT est associée la coupure des rationnels qui lui sont inférieurs ; la complétude de la réalisation ordonnée fournit l’unique élément ayant cette coupure. La densité rationnelle prouve l’injectivité et la conservation stricte de l’ordre. Le même procédé en sens inverse prouve la surjectivité. Les approximations rationnelles par excès et par défaut conservent la somme et le produit, et donc la structure de corps. L’élément positif dont le carré est \(a>0\) est transporté vers l’élément positif dont le carré est \(\iota(a)\). Si une suite converge, pour chaque tolérance rationnelle positive ses derniers termes sont dans le voisinage rationnel de la limite ; la conservation de l’ordre transporte cette condition. Toutes ces opérations s’effectuent dans \(\mathfrak G\). L’unicité découle de la conservation des coupures rationnelles.
\end{proof}

\subsection{Descente par les fibres et conservation du sélecteur}
\label{hmtfg:descenso-selector}

% XI REV11, ch_descenso_por_fibras.tex:10--39; prueba completa pertinente.
\begin{proposition}[Critère de descente d’une opération]
\label{hmtfg:criterio-descenso}
Soient \(q:\Omega\twoheadrightarrow J\) et
\(\star:D\subseteq\Omega^2\to\Omega\), avec \(D\) saturé par les fibres de \(q\times q\). Il existe une unique opération \(\bar\star\) sur \((q\times q)(D)\) telle que
\[
 q(x\star y)=q(x)\,\bar\star\,q(y)
\]
si et seulement si
\[
 q(x)=q(x'),\quad q(y)=q(y')
 \ \Longrightarrow\ q(x\star y)=q(x'\star y')
\]
pour les couples du domaine.
\end{proposition}
\begin{proof}
Si \(\bar\star\) existe, les deux résultats sont l’évaluation du même couple d’arguments. Réciproquement, on choisit des antécédents des arguments visibles et on définit le résultat comme la lecture de leur opération. La saturation fixe le domaine ; la condition sur les fibres rend le résultat indépendant des antécédents choisis. La surjectivité prouve l’unicité.
\end{proof}

Dans une application au moyen de \(P_{\rm ar}\), ce critère identifie la condition précise permettant de transporter une opération depuis les états complets jusqu’à leur lecture. Dans la carte \(\iota\), la structure de corps et les limites satisfont déjà cette condition. Le profil analytique est ensuite construit à partir du calendrier, et cette compatibilité est utilisée pour transporter son itération et ses racines.

\subsection{Lecture dodécaphasique et normalisation quadratique}
\label{hmtfg:perfil-dodecafasico}

% Propietario integral c37_feigenbaum.tex: comienzo hasta la familia f_lambda.
La section orientée de criticité sur les douze fenêtres du calendrier est
\[
 \mathcal P_{\rm crit}(m)=\varphi_m
   =(30m-10)\frac{\pi_{\rm HMT}}{180}
   =\pi_{\rm HMT}\frac{3m-1}{18},\qquad m=1,\ldots,12.
\]
Cette section fixe l’origine et l’orientation de ses représentants angulaires. Le lecteur uniforme des douze fenêtres est
\[
 \eta_0(F)=\frac1{12}\sum_{m=1}^{12}F(\varphi_m).
\]
On distingue la fonction \(\eta_0\), qui attribue des poids à ces fenêtres, de la carte ordonnée \(\iota\). La normalisation uniforme et la section angulaire font partie des données operatorielles de ce profil.

En posant \(a_m=(3m-1)/18\), on obtient
\[
 D_0=\frac1{12}\sum_{m=1}^{12}a_m^2=\frac{899}{648},
 \qquad
 s_0=\frac{36}{\pi_{\rm HMT}\sqrt{899}},\qquad
 k_m=s_0\varphi_m=\frac{2(3m-1)}{\sqrt{899}}.
\]
En effet, \(\sum_{m=1}^{12}(3m-1)^2=5394=6\cdot899\), de sorte que \(s_0^2\pi_{\rm HMT}^2D_0=2\). Le profil est
\begin{equation}
 u_0(x)=\frac1{12}\sum_{m=1}^{12}\bigl[1-\cos(k_mx)\bigr],
 \qquad f_\lambda(x)=1-\lambda u_0(x).
 \label{hmtfg:perfil-familia}
\end{equation}
La constante \(\pi_{\rm HMT}\) intervient comme lecture angulaire du noyau HMT et s’annule algébriquement dans la normalisation du profil. L’échelle \(\lambda\) est le paramètre de la famille, dont la valeur critique sera sélectionnée par la dynamique.

\begin{proposition}[Normalisation et forme fermée du profil]
\label{hmtfg:forma-perfil}
Le profil \(u_0\) est pair et analytique, satisfait
\(u_0(0)=u_0'(0)=0\) et \(u_0''(0)=2\), et, pour
\(y=x/\sqrt{899}\), admet l’expression
\[
 u_0(x)=1-\frac{\cos(37y)\sin(36y)}{12\sin(3y)}
       =1-\frac{\sin(73y)-\sin(y)}{24\sin(3y)}.
\]
Les singularités apparentes s’interprètent par continuité. Son développement initial est
\[
 u_0(x)=x^2-\frac{715769}{2424603}x^4+O(x^6).
\]
\end{proposition}
\begin{proof}
La somme finie de cosinus conserve la parité et l’analyticité. Sa dérivée seconde en zéro est \(12^{-1}\sum_m k_m^2=2\). La somme de la progression trigonométrique \(\sum_{m=1}^{12}\cos((6m-2)y)\) est \(\cos(37y)\sin(36y)/\sin(3y)\) ; l’identité produit--somme fournit la seconde expression. La série du cosinus et \(\sum_m(3m-1)^4=4294614\) donnent le coefficient de degré quatre indiqué. L’expression comme somme finie fixe le prolongement aux zéros du dénominateur de la forme fermée.
\end{proof}

\subsection{Transport des itérées, des racines et des minima}
\label{hmtfg:transporte-raices}

% Reunión existente: TRANSVERSALIDAD_Y_ESCALADO_LOCAL_FEIGENBAUM.md, sección 14.1.
Dans \(\mathbb R_{\rm HMT}^{\mathfrak G}\), la série convergente
\[
 \cos_{\rm HMT}(t)=\sum_{j=0}^{\infty}
                   \frac{(-1)^j t^{2j}}{(2j)!}
\]
définit la réalisation analytique correspondante. Soient
\[
 u_{\rm HMT}(x)=\frac1{12}\sum_{m=1}^{12}
  \left[1-\cos_{\rm HMT}
   \left(\frac{2(3m-1)x}{\sqrt{899}_{\rm HMT}}\right)\right],
 \qquad F_a(x)=1-a\,u_{\rm HMT}(x),
\]
\[
 S_n^{\rm HMT}(a)=F_a^{\,2^n}(0),\qquad
 S_n(\lambda)=f_\lambda^{\,2^n}(0).
\]

\begin{proposition}[Conjugaison ordonnée de la famille et du sélecteur de racines]
\label{hmtfg:conjugacion-selector}
Pour les arguments internes des expressions précédentes,
\[
 \iota\circ F_a=f_{\iota(a)}\circ\iota,\qquad
 \iota(S_n^{\rm HMT}(a))=S_n(\iota(a)).
\]
La carte transporte exactement les conditions de retour, de période exacte et d’ordre des paramètres. Si un ensemble de racines admissibles possède un minimum, l’image de ce minimum est le minimum de l’ensemble image.
\end{proposition}
\begin{proof}
La carte fixe les rationnels, transporte la racine carrée positive et conserve les limites des sommes partielles du cosinus. Ainsi, \(\iota(u_{\rm HMT}(x))=u_0(\iota(x))\). La compatibilité avec la somme et le produit prouve la première égalité ; la récurrence sur le nombre d’itérations prouve la seconde. Une racine se caractérise par l’égalité à zéro, conservée dans les deux sens. La période exacte \(2^n\) ajoute les inégalités \(F_a^{2^j}(0)\ne0\), \(0\le j<n\), également conservées par l’injectivité de \(\iota\). Si l’on exige que le paramètre dépasse un centre précédent, cette inégalité est transportée par conservation de l’ordre. Enfin, \(a\le b\) pour tout \(b\) de l’ensemble équivaut à \(\iota(a)\le\iota(b)\) pour tout élément de son image. La surjectivité sur les ensembles ainsi définis complète l’affirmation relative au minimum.
\end{proof}

Le domaine de ce transport est le continu interne indiqué. L’existence de chaque minimum et son identification à une branche dynamique précise sont des résultats de la construction de criticité qui suit ; la carte conserve ces propriétés lorsqu’elles ont été établies.

\subsection{Représentation spectrale finie du lecteur uniforme}
\label{hmtfg:jacobi-finito}

% Grassmannianos REV07, sections/jacobi_spectral.tex:35--136,
% mecanismo de Hankel, Gram--Schmidt, recurrencia y representación espectral.
% Especialización atómica ya reunida en la sección 14.2 del desarrollo Feigenbaum.
Les fréquences au carré du profil définissent la mesure positive
\[
 s_m=k_m^2=\frac{4(3m-1)^2}{899},\qquad
 \nu_0=\frac1{12}\sum_{m=1}^{12}\delta_{s_m},\qquad
 M_r=\int s^r\,d\nu_0(s).
\]
Ses douze nœuds sont distincts et positifs. Pour \(1\le r\le12\), la matrice de Hankel \(H_r=(M_{i+j})_{0\le i,j<r}\) est définie positive. En effet, pour \(p\) de degré inférieur à \(r\),
\[
 \mathbf p^{\mathsf T}H_r\mathbf p
   =\frac1{12}\sum_{m=1}^{12}p(s_m)^2>0
\]
sauf si \(p=0\), car un polynôme de degré au plus onze qui s’annule en douze nœuds distincts est nul. Au degré douze apparaît le polynôme \(\prod_m(s-s_m)\), de norme nulle pour cette mesure. La représentation spectrale correspondante est de dimension douze.

Soient \(p_0,\ldots,p_{11}\) les polynômes orthogonaux moniques,
\(h_j=\int p_j^2\,d\nu_0>0\) et
\(\psi_j=p_j/\sqrt{h_j}\). Comme \(\nu_0\) est une probabilité,
\(\psi_0=1\). La multiplication par \(s\) dans \(L^2(\nu_0)\) est autoadjointe et possède une matrice tridiagonale \(J_{12}\) dans cette base.

\begin{proposition}[Récurrence et entrelacement spectral exacts]
\label{hmtfg:jacobi-entrelazamiento}
Avec
\[
 a_j=\frac{\int s p_j(s)^2\,d\nu_0(s)}{h_j},\qquad
 \beta_j=\frac{h_j}{h_{j-1}}>0\quad(1\le j\le11),
\]
la matrice \(J_{12}\) a pour diagonale \(a_j\) et pour sous-diagonale
\(\sqrt{\beta_j}\). Si
\[
 Q_{mj}=\frac{\psi_j(s_m)}{\sqrt{12}}
 \quad(1\le m\le12,\ 0\le j\le11),\qquad
 D=\operatorname{diag}(s_1,\ldots,s_{12}),
\]
alors
\[
 Q^{\mathsf T}Q=I,\qquad DQ=QJ_{12},\qquad
 J_{12}=Q^{\mathsf T}DQ,
\]
et le profil satisfait
\begin{equation}
 u_0(x)=e_0^{\mathsf T}
          \bigl[I-\cos(x\sqrt{J_{12}})\bigr]e_0.
 \label{hmtfg:perfil-espectral}
\end{equation}
\end{proposition}
\begin{proof}
La positivité des systèmes de Hankel fournit de manière unique les polynômes orthogonaux moniques. En développant \(s p_j\) dans la base orthogonale, les coefficients de \(p_i\), \(i\le j-2\), s’annulent car
\(\langle sp_j,p_i\rangle=\langle p_j,sp_i\rangle=0\).
Le coefficient de \(p_{j+1}\) est un par monicité, celui de \(p_j\) est \(a_j\) et celui de \(p_{j-1}\) est \(h_j/h_{j-1}\). Au dernier degré, la même égalité s’entend dans \(L^2(\nu_0)\), où le polynôme monique de degré douze qui s’annule aux nœuds représente le vecteur nul. La normalisation donne la sous-diagonale positive indiquée.

L’orthonormalité équivaut à \(Q^{\mathsf T}Q=I\). L’égalité de récurrence à chaque nœud donne \(DQ=QJ_{12}\). Comme \(Q\) est carrée, on a également \(QQ^{\mathsf T}=I\), et l’on obtient la diagonalisation de \(J_{12}\). Ses valeurs propres positives permettent de définir sa racine carrée positive. Pour toute fonction définie en ces douze nœuds,
\[
 e_0^{\mathsf T}\Phi(J_{12})e_0
   =\frac1{12}\sum_{m=1}^{12}\Phi(s_m),
\]
car \(Qe_0=(1,\ldots,1)^{\mathsf T}/\sqrt{12}\).
Le choix \(\Phi(s)=1-\cos(x\sqrt{s})\) prouve
\eqref{hmtfg:perfil-espectral}.
\end{proof}

La représentation conserve exactement la fonctionnelle uniforme, ses moments et son profil. Ses premiers coefficients sont
\[
 a_0=M_1=2,\qquad
 \beta_1=M_2-M_1^2=\frac{2493348}{808201}.
\]
La généalogie complète demeure dans l’état TPK dont provient le lecteur ; \(J_{12}\) représente sa lecture spectrale particulière. La construction de Gram--Schmidt et la récurrence de Jacobi s’appliquent à cette mesure atomique, avec leur terminaison en dimension douze.

La chaîne operatorielle utilisée dans les sections suivantes est ainsi fixée : les cellules APP et l’orientation TRIT alimentent la transition TPK ; ses prolongements produisent les cylindres et leur lecture ordonnée ; les douze fenêtres orientées et la fonctionnelle uniforme produisent \(u_0\) ; la famille \(f_\lambda\) détermine les retours et les racines ; la carte interne conserve l’ordre de leurs sélecteurs. La réalisation spectrale décrit le même profil au moyen d’une matrice autoadjointe finie et d’un vecteur cyclique.
